% !TEX program = xelatex
\documentclass[UTF8,a5paper,11pt,openany]{ctexbook}

% === Page Layout ===
\usepackage[paper=a5paper,top=2.2cm,bottom=2.2cm,left=2cm,right=2cm]{geometry}

% === Fonts ===
\setCJKmainfont{SimSun}[BoldFont=SimHei,ItalicFont=KaiTi]
\setCJKsansfont{SimHei}
\setCJKmonofont{FangSong}
\setmainfont{Times New Roman}

% === Headers and Footers ===
\usepackage{fancyhdr}
\pagestyle{fancy}
\fancyhf{}
\fancyhead[LE,RO]{\small\ttfamily\leftmark}
\fancyfoot[CE,CO]{\small\thepage}
\renewcommand{\headrulewidth}{0.4pt}
\fancypagestyle{plain}{
    \fancyhf{}
    \fancyfoot[CE,CO]{\small\thepage}
    \renewcommand{\headrulewidth}{0pt}
}

% === Chapter Styling ===
\usepackage{titlesec}
% Fix first paragraph indent after chapter headings (titlesec overrides ctex default)
\usepackage{indentfirst}
\usepackage{etoolbox}
\makeatletter
\patchcmd{\@afterheading}{\@afterindentfalse}{\@afterindenttrue}{}{}
\makeatother

\titleformat{\chapter}[hang]
    {\LARGE\bfseries\heiti}
    {\chaptertitlename\ \thechapter}{1em}{}
\titlespacing*{\chapter}{0pt}{-30pt}{20pt}

% === Scene Break ===
\newcommand{\scene}{\vspace{1em}\begin{center}*\quad*\quad*\end{center}\vspace{0.5em}}

% === Content Depth in TOC ===
\setcounter{tocdepth}{0}

% === Hypersetup ===
\usepackage{hyperref}
\hypersetup{
    colorlinks=true,
    linkcolor=black,
    filecolor=black,
    urlcolor=black,
    citecolor=black,
    pdfauthor={},
    pdftitle={下沉},
    pdfsubject={当代现实小说},
    pdfkeywords={下沉, 非对称性, 当代小说},
    bookmarksnumbered=true,
    bookmarksopen=true,
    pdfencoding=unicode,
}

% === Spacing ===
\setlength{\parindent}{2em}
\setlength{\parskip}{0em}

\begin{document}

% === Title Page ===
\begin{titlepage}
\vspace*{4cm}
\begin{center}
    {\fontsize{36}{44}\bfseries\heiti 下\quad 沉\\[0.5cm]}
    \vspace{2cm}
    {\Large\kaishu 一条河。一个渡口。十年。}\\[1cm]
    \vspace{3cm}
    {\normalsize\kaishu 方远 著}\\[0.3cm]
\end{center}
\end{titlepage}

% === Table of Contents ===
\cleardoublepage
\tableofcontents

\cleardoublepage


% ==== Chapter 1 ====
\chapter{辞职}
\indent
办公室里没多少东西。

方远拉开左边的抽屉。一把钢尺，用了六年，刻度磨掉了一半。一支钢笔，笔帽上的漆落了一小块。一叠便签纸，还剩七八张。他把便签纸放进纸袋里。

窗外是江。

江水的颜色随天气变。晴天泛白，阴天发灰。今天是阴天，江水是灰的。

他在规划局坐了六年。分管工业用地那一块。桌上堆过图纸、申请、会议纪要。他说过不少话，但此刻想不起来的居多。倒是有几张图纸上的线条还在脑子里。土地怎么划，路怎么走，厂建在哪。这些东西看多了，闭上眼睛也能画出来。

他把钢尺放进纸袋。

走廊上有人在复印。机器嗡嗡响了一阵，停了，又响。那个声音他听了六年，今天是最后一次。

钢笔没墨了。他在便签纸上划了两道，什么都没有。把笔也放进纸袋里。

桌上空了。键盘、电话、座签，都不是他的东西。

他去窗边站了一会儿。

江水往下游流。远处有几艘铁壳船，看不清是停在岸边还是在动。对岸是开发区，烟囱不冒烟已经两年了。有一片新楼盘盖了一半，塔吊立着，没转。下面有几只白色的鸟绕着塔吊飞。飞了很久也没有落下来。

六年前他来这里报到的时候，那片开发区还是图纸上的一块红框。

办事员小吴推门进来，问他需不需要帮忙。他说不用。小吴站了一会儿，没说话，出去了。

他把抽屉依次拉出来检查了一遍。第三格抽屉里有一张照片。他在办公室楼下拍的。那年刚入职，穿了一件不合身的夹克，站在楼前的花坛旁边。花坛现在改成停车场了。

他看着那张照片。照片里的人没什么表情，不太像在拍照，更像在办什么手续。他看了几秒。把照片放进纸袋里。

走廊上的复印机停了。

他提着纸袋走出去。电梯来了，他看了一眼，走了楼梯。

一共五层。每一层的楼梯间有窗户，窗外是同一段江面。他一层一层地往下走。走到二楼的时候，停了一下。二楼是审批科。门关着。里面有人说话。他继续往下走。

一楼大厅的保安认识他。保安正在看手机，抬头冲他点了点头。他也点了点头。

推开玻璃门，秋天扑面而来。


\scene

鲁建国等在门口。

风很大。鲁建国缩着脖子，手里夹着一根没点的烟。看到方远就收了烟，说，走吧。

方远说，嗯。

两个人沿着江边的路走。梧桐叶落了一地，踩上去有细碎的响声。江水在左手边，灰蒙蒙的，没有船。偶尔有电动车从身后经过，按一声喇叭，就不见了。

走了几分钟，鲁建国说，想吃面。

方远说，好。

面馆离规划局不远。在一条巷子里，门脸很窄，招牌褪了色。陈敏系着围裙站在灶前，看见他们进来，点了点头。

坐吧。她说。

店里没别的客人。头顶的吊扇转得很慢，发出轻微的吱嘎声。墙上贴着几张招贴画，边角卷了。有一张是去年过年的年画，还没摘下来。娃娃抱鱼的图案褪成了淡粉色。

陈敏端来两杯水。水是温的，带着玻璃杯用久了特有的味道。

鲁建国说，素面，多加辣椒。

方远说，一样。

陈敏进厨房了。灶火的声音传出来，锅铲碰铁锅，水开了倒进锅里。这些声音方远很熟。六年来听了无数次，但今天和以前不一样。以前只是声音。今天他注意到锅铲碰锅的那一声，在灶火响起来之前有一个很短的间隙。那个间隙里，厨房是完全安静的。

鲁建国把烟掏出来，看了看墙上贴的禁烟标牌，又放回去。

他问，东西都拿完了。

方远说，拿完了。

鲁建国说，六年就这点东西。

方远说，嗯。

面来了。热气扑在脸上。陈敏放下碗的时候，多看了一眼方远。他没看她。她在围裙上擦了一下手，回厨房去了。

两个人都没说话。筷子起落，面一根一根地减少。鲁建国吃辣椒吃得额头出了汗，用袖子擦了一下。方远吃得慢一些。汤底有酸菜。不是很多，刚好够味。

吃得差不多了，鲁建国放下筷子。他看了一眼窗外。窗外没什么好看的。一只猫蹲在对面的墙头上。猫的耳朵动了一下，警惕着某个声响，然后继续趴下去。

他说，你接下来打算怎么办。

方远没抬头。把最后一口汤喝了。

汤有点咸。

他说，走走看看。

鲁建国没接话。过了一会儿，把剩下的面吃完了。他的碗比方远干净。汤底一滴不剩。

方远去结账。陈敏收钱的时候看了他一眼。他在这家面馆吃了六年，每星期至少两次。她什么都不问。只是找钱的时候，多给了他一张纸巾。

方远接过纸巾。不是那种面馆里常备的劣质纸巾，是带一点花纹的那种。他把纸巾握在手里，没擦嘴。

说了声，谢谢。

陈敏说，慢走。


\scene

出了面馆，天已经黑了。

路灯亮了。黄色的光落在江面上，拉出细碎的波纹。风比下午小了些。远处有人在放音乐，声音很远，听不清是什么歌，只能辨认出节奏。

鲁建国说，那我先回了。

方远说，好。

鲁建国走了两步，回头看了一下。方远站在路灯下，提着纸袋。鲁建国说了一句什么，声音不大，风一吹就散了。

方远没听清。也没问。

他沿着江边走。

这条路走了六年。早晨来上班，傍晚回去。春天有卖菜的推着车经过。夏天江边的石栏被晒得发烫。秋天梧桐叶铺了一地。冬天江风硬得刺骨。

前后左右，都是一样的。

不一样的只有江水。江水每天都不一样。今天的水位比昨天低了半米。明天会更高还是更低，不知道。

他在石栏边停下来。

路过的人从他身后走过。有抱孩子的，有牵着狗的，有边走边打电话的。他对这些声音充耳不闻。江水的声音不大，但一直都在。

对岸的开发区亮了几盏灯。不是很多。五盏。也许六盏。有一盏忽明忽暗，可能是接触不良。

他在规划局经手过这个开发区的审批。最早的那批企业，有一半搬走了。剩下的那些，有的在生产，有的只是占着地块等涨价。涨没涨他也不知道。他已经不在那里了。

他往上游看去。

江面在黑暗中延展。看不出多远。大约过了那片新建的水闸，再往上就是老工业区了。他记得那里的化工厂、轴承厂、印染厂。大部分已经停产。厂房还在。

水还在流。

他站了一会儿。风从下游吹来，带着一种说不清楚的味道。像铁锈，又像旧木头泡了水。

纸袋不重。六年攒下来的东西，一个纸袋就装完了。

他把纸袋换到左手。右手插进外套口袋里，摸到一包烟。不是他的。

鲁建国留下的。

他抽出一根。没有火。

他把烟放在江风的出口处，看着它被风吹动，一点点旋转。然后松开了手指。

烟落进水里。没有响动。

水往东流。

他转身往回走。身后是路灯，眼前是黑的路。走了几十米，路灯的光越来越远，最后只剩江面的反光。那种光不是亮，是一种比黑暗浅一点的灰。靠着它，勉强能看到路。

他在黑暗中走了一段。

然后上了大路。


\scene

家里的灯是他早晨走的时候开的。开了整整一天。

他换了鞋，把纸袋放在桌上。桌上还有一份没有收的报纸。昨天的。头版是开发区的新规划。又一轮调整。他没看。

他去厨房烧了一壶水。等水开的时候，靠在灶台上。

灶台的瓷砖有一块裂了。是夏天的时候裂的。他没修。裂缝里积了灰，灰是灰白色的。

水烧开了。电水壶自动跳了。他泡了一杯茶。

茶叶是陈的，在杯底泡了很久才散开。

他端着茶走到窗前。这扇窗看不到江。窗外是一堵墙。墙的那边是另一栋楼。楼里有人在做菜。油锅的声音。

他喝了口茶。很烫。

手机震了一下。

鲁建国发来的消息。两个字。

保重。

他看着这两个字。鲁建国不是会发这种话的人。这种话对他来说太正式了。可能他也觉得今天不一样。可能他也觉得应该说点什么。但想了很久，也只有这两个字。

方远把手机翻过去，屏幕朝下放在桌上。

茶凉了。

他没有再添水。


\scene

第二天他起得很早。

天还没亮透。窗外灰蒙蒙的，分不清是雾还是天光。他套了一件旧棉袄，锁了门。

楼下早点摊刚支起来。蒸笼冒着白汽。老板蹲在地上点火，看他过来，有点意外。

今天早。

他嗯了一声。买了两根油条，一杯豆浆。豆浆烫手。他站在路边喝完，把塑料杯扔进垃圾桶。油条吃了一根半。剩下半根拿在手里，凉了。

路边等了五分钟。第一班公交车来了。

车上没人。他挑了最后一排靠窗的位置。

车发动的时候，整个车身抖了一下。然后慢慢开出城区。

穿过老城区的时候，有早起遛狗的人牵着狗过马路。司机按了喇叭。狗回头看了一眼，没让。司机又按了一下。遛狗的人这才拽了拽绳子。

出了城就是江。

江水比昨天清了些。也许是天光的问题。早晨的江面泛着一层薄薄的青色。不像是脏，像是还没醒。

他在一个没有名字的站下了车。站牌歪了。上面的字被风雨磨掉了一半，能认出来的只有一个江字。

路是土路。两边长着过膝的草。草地上有野花的残枝。秋天快结束了，花都枯了。枯花锈成一团，风一吹就散了。有一股枯草烧过的气味。

他沿着路往下游走。

大约走了四十分钟，听到了另一种水声。不是江的水声。江的水声是沉闷的、持续的、像呼吸。这个声音是碎的、急的，是水撞在什么东西上面。

他穿过一片芦苇。

面前是一个渡口。

水泥码头裂了缝。缝里长着草。有三条锈蚀的缆桩，缆绳早没了。旁边有一间铁皮房子，门虚掩着。房子后面堆着几根腐烂的木桩，上面长了青苔。

渡口过去应该是热闹的。码头上有车辙的痕迹，很深，被水泥封住了。但现在已经没人用了。上下游都修了新桥。桥通了，渡口就废了。

但这个位置很好。

方远站在码头上，往下游看。从这个角度看过去，江的走势一清二楚。上游几个工厂的位置，下游那个规划中的物流园的位置，全都在视线之内。

他看了一会儿。

然后掏出手机。打开地图。

他放大了江城这一段。那个物流园，在规划图上只标了一个虚线框。但他知道规格。他当初参与了论证会。货物吞吐量、日均车辆数、仓储面积，他都有数。不是记在纸上，是记在脑子里。

渡口的位置，在物流园上游三公里。

三公里，刚好是一个不需要额外审批的距离。

他把地图关掉了。

江风从下游吹来。芦苇沙沙响。那种声音不好形容。不像音乐，不像噪音。像很多人在远处小声说话。你听不懂他们在说什么，但你知道他们在说。

他在码头上坐下来。水泥很冷。冷透过裤子，传到皮肤上，然后就没有然后了。人的身体会适应冷。

他看着江水。看了大约二十分钟。

然后站起来。拍了拍裤子上的灰。

往回走。

回去的路上他走得很慢。有时候停下来，看看岸边的水位线。水位线有好几层。最高的一层在石壁上，上面长了青苔，颜色比其他地方深。有时候蹲下来，摸一摸土壤。土是湿的。不是雨打湿的。是江水渗上来的。

快到公交站的时候，他看见路边一个老头在卖菜。地上铺了一块塑料布，摆着几把青菜、几个萝卜。青菜的叶子有点蔫。老头蹲在旁边，手里夹着一根烟，烟雾和早晨的雾气混在一起。

他买了一把青菜。

老头找不开零钱。他说不用找了。

上了公交车。还是同一个位置。最后一排，靠窗。

江水在窗外流着。他闭上了眼睛。手里握着那把青菜。菜叶从塑料袋里露出来，垂在他的膝盖上。


\scene

鲁建国打来电话，是第四天。

他说，你在哪。

方远说，在江边。

鲁建国沉默了一下。说，你那天晚上说的走走看看，是什么意思。

方远说，就是走走看看。

电话那头又是沉默。不是那种不知道该说什么的沉默。是那种话已经想好了，但还没决定要不要说的沉默。

方远说，你来一下。

鲁建国说，去哪。

方远说了那个渡口的位置。

两个小时后，鲁建国来了。他开着一辆旧的黑色轿车，轮胎上沾满了泥。下车的时候踩了一脚水，骂了一句。

方远站在码头上。鲁建国走过去，和他并排站着。

鲁建国说，这是什么地方。

方远说，渡口。

鲁建国看了看四周。芦苇、铁皮房子、锈了的缆桩。江风吹得芦苇弯了腰。他说，我知道是渡口。你带我来这干什么。

方远没回答。他指了指下游。

说，那边三公里，物流园。

鲁建国顺着他的手指看过去。什么都没看到。江面和对岸的雾气混在一起，分不清界限。他眯着眼睛看了很久，最后放弃了。

鲁建国说，然后呢。

方远说，新桥通了，渡口废了。但水没变。河道没变。货物流向没变。路绕远了，水是直的。

鲁建国看着方远。看了好一会儿。

他说，你要搞货运。

方远说，嗯。

鲁建国又想了一下。说，你哪来的钱。

方远说，存了六年。

又想了想。

鲁建国说，你不是走走看看。

方远没说话。

江面上有一只水鸟。白色的，飞得很低，贴着水面。翅膀尖几乎蹭到水。它飞了一会儿，翅膀扇动的频率很慢，慢得像是在滑行。然后停在一根露出水面的树桩上。树桩上站着另外两只鸟。第三只落下来的时候，有一只飞走了。

鲁建国说，我老婆不会同意的。

方远说，我没让你现在来。

鲁建国说，我现在也不会来。

又停了。

鲁建国从口袋里掏出烟。这次点了。吸了一口。烟从鼻子里出来，被风吹散了。他吸得不快，像是要把每口烟都留在肺里多一点时间。

鲁建国说，但我觉得你会做成。

方远笑了笑。那个笑容很短。嘴角动了一下就收回去了。鲁建国甚至不确定自己是不是看到了。他没有笑出声。方远笑的时候不出声。

方远说，你要是哪天想来了，跟我说一声。

鲁建国说，好。

他把烟抽完。烟蒂在鞋底捻灭，放进自己的口袋里。地上什么都没有留下。

两个人沿着江边走回车里。

路上谁都没再说话。

车发动了。老旧引擎的声音闷在引擎盖里，突突地响。鲁建国调了个头，泥水溅起来，落在轮胎上。

他摇下车窗。

方远。

嗯。

那天你在路灯下说了句什么话。风太大我没听清。

方远想了想。

他说，我说的是：水还在流。

鲁建国看着他。他坐在驾驶座上，没说话，只是看着方远。那个眼神里有些东西。方远看得出来，但他没接。

然后鲁建国升起车窗。车子开走了。

尾灯在土路上颠簸了两下，拐过一个弯，看不见了。

方远站在原地。

风停了片刻。很短的一瞬。芦苇也不响了。然后风又开始吹。河水的声音回来了。那种沉闷的、持续的、像呼吸一样的声音。它从来没停过。只是有时候风盖过它，让人觉得它停了。

他转身走向渡口。

铁皮房子的门还是虚掩着。他推开门。屋里黑洞洞的。一间小屋子，一张破桌子，一个生锈的铁柜。地上有几只死虫子。墙上挂着一九五几年的安全守则，纸变成了褐色，字一个都看不清。

他从兜里掏出手电筒。

照了一圈。

墙角有一只老鼠窜过去。很快。只听到声音，没看清。手电筒的光追上它的时候，它已经不见了。

他没有动。

他想起他的父亲。父亲在船上工作。有时候带他上船。他坐在船舱里，听到水拍打船帮的声音。那种声音和现在听到的一样。

父亲说，河不会骗人。

父亲还说过别的话。他在那个铁皮房子里站了一会儿，试图想起来。但只记得这一句。

河不会骗人。

他关掉了手电筒。站在黑暗里。河水的声音从房子的铁皮墙外面渗进来。不是很大，但很完整。

他站了一会儿。

然后推门出去了。

\clearpage

% ==== Chapter 2 ====
\chapter{渡口}
\indent
铁皮房子的门关了。

方远站在码头上。阳光在水面上跳。十一月的太阳不高，光是从侧面过来的，把芦苇的影子拉得很长。

他需要知道这块地是谁的。

在规划局六年，他知道怎么查。不是那种官方的查询方式。是更快的。打过交道的，吃过饭的，欠过人情的。每个办公室里都有这样的人。你给他打过一次电话，他就欠了你一份人情。你不急着要他还。

他翻出手机里的通讯录。往下滑了很久。

有时候号码多了，找一个人和丢一个人一样难。同音字多了，滑着滑着就忘了自己在找谁。

他拨了一个号码。响了五声。没人接。

他又拨了一次。这次响了两声。

喂。一个沙哑的声音。

老唐。是我。方远。

谁。

方远。规划局的。

哦。方工。好久没听到你声音了。什么事。

想打听一块地。

哪块。

方远说了那个渡口的位置。电话那头沉默了一下。老唐在脑子里翻找。方远听得到对方的呼吸声。

那个啊。老唐的声音慢下来，那一片前年划给水利局下属的一个单位了。不过他们没开发。账上没钱。

产权还在他们手里吗。

应该在。你要干什么。

想用。

方远听见老唐那边有翻纸张的声音。老唐在找东西，但可能只是习惯。他的桌上永远堆着半尺高的文件。

那一片不是什么好地方啊。老唐说。

嗯。

路也不通。

嗯。

你要用多久。老唐问。

方远说，能多久就多久。

老唐没再问。干了二十多年，他知道什么时候不该问。

我帮你打听一下。三天。老唐说。

好。

方远挂了电话。把手机放进口袋。

他坐在码头上等。不是等电话。只是不知道该去哪。家里空空荡荡，单位回不去了。只剩下这里。

江水在他面前流。不急，也不慢。上游漂下来一根断枝。在码头边转了两圈，又被水带走了。


\scene

老唐回电话是在第三天晚上。

方远正在烧水。电水壶嗡嗡响。电话响了。他等水烧开了，泡了茶，才接。

打听到了。老唐说，他们愿意租。一年。

方远说，一年不够。

我知道。老唐说，但他们只敢签一年。上面查得紧。你租一年，做起来了，再续。

方远喝了口茶。

租金多少。

他们开价一年八万。老唐笑了一声。我压到了五万。但他们加了个条件。合同期满如果续租，你有优先权。但如果有人出价比你高，你得跟着加。

方远说，可以。

还有一件事。老唐停顿了一下。他们要求签三方协议。除了你和水务单位，还需要区交通局出具一份意见函。说你用的这块地不占用航道。

方远的手指在茶杯上停了一下。

交通局。

对。老唐说，你以前跟他们打过交道吗。

方远说，打过。

那就行。

老唐报了对方的联系方式。一个座机号码，一个名字。方远没记在纸上，只是听了一遍。他听的时候说了两遍，就记住了。然后把号码存进了手机。

老唐说，还有别的事吗。

没了。谢谢。

不客气。老唐说。然后加了一句：那个渡口，我年轻的时候还坐过船。船是木头的。柴油机的声音很大。

方远说，现在没有船了。

是啊。老唐说。挂了。

方远把手机放在桌上。屏幕暗了。

水又流了三天。他在等别的东西。


\scene

一个星期后，钱到账了。

他从账户里转出五万块。那是他六年存下来的一部分。不算多，也不算少。刚好够开始。

签合同那天，他穿了一件干净的白衬衫。但那件衬衫在衣橱里压了太久，有褶子。他没有熨斗。他用手抹了两下，没用。

水务单位的办公室在老城区一栋灰楼里。二楼。楼梯间堆着旧报纸和破椅子。空气里有潮湿的味道。不是霉，是旧。所有东西凑在一起散发出的那种旧。

接待他的是一个中年女人。戴眼镜。说话不紧不慢。

方远坐在她对面的木椅上。椅子硬，坐得不舒服。

她看了他的材料。抬头看了他一眼，又低头看材料。

她说，你之前是规划局的。

嗯。

那个铁皮房子，你看过了吗。

看过了。

里面什么都没有。

嗯。

她没再说下去。把合同推过来。

合同不长。两页纸。用词很古老。里面有一个词叫宁可。他有很多年没在正式文件里见过这个词了。

他签了字。钢笔在纸上划过的声音很轻，但在安静的办公室里显得很大。

她盖了章。红色的。

他拿到了一份复印件。复印机卡纸了。中年女人弯腰去修，拍打了机器好几下。方远站起来想帮忙，她已经修好了。

她说，预祝顺利。

方远点了点头。

走出那栋灰楼，他在外面的台阶上站了一会儿。阳光刺眼。旁边有个卖烤红薯的推车，煤炉里的炭火发出暗红色的光。他买了一个。红薯烫手，他在两只手之间来回倒着，边走边吃。

红薯的皮烤焦了，里面的肉是黄的，甜得发腻。

他吃完了。把皮扔进路边的垃圾桶。

那一刻他想到了铁皮房子。想象它亮着灯的样子。想象里面有一张好桌子。想象码头上停着船。

这个想象很短，大概只持续了几秒钟。

然后就没了。

他继续走。


\scene

开张前的第一件事，是收拾。

他买了一辆二手面包车。花了八千。车身上有几处刮痕，但发动机还行。他把工具塞进后备箱。铁锹、扫帚、水桶、抹布、铁锤、钉子。还有一个旧的柴油发电机。卖家说这个发电机用了不到两年，声音有点大，但能用。他信了。

到了码头是早上七点。雾还没散。江面上白茫茫的，看不见对岸。

他先修了铁皮房子的门。门铰链锈了，开门时整个门往下坠。他用铁锤敲松了锈死的螺丝，换上了新铰链。试了一下。门开了，不会坠。

然后是屋子里。死虫子扫出去。生了三簸箕。他没戴口罩。灰尘扬起来，呛得他咳了几声。旧桌子搬到外面，用河水冲洗了桌面。桌面暗红，上面的水渍像年轮一样一圈一圈散开。铁柜子也一样，冲洗后放在太阳底下晒干。

码头上的裂缝他没管。太重了。他需要水泥和沙子。

他从路边捡来石块，把码头边缘碎掉的地方垫平了。石块大小不一，拼在一起像是补丁。

然后他站远一些看。

还是破的。

但没那么破了。

路边的草长到了码头边缘，他弯腰拔了一阵。拔到一半，手疼。他停下，直起腰，看了一会儿江面。又继续拔。

下午的时候，他把铁皮房子通了电。

从最近的电线杆上接下来，用了不到两百米的电缆。接电的手续他走了三天。报到供电所，填了表，等通知。通知来了，他就拉了。

灯亮了。

是一盏很普通的白炽灯。四十瓦。灯光黄黄的。他站在灯光下，看着自己的影子落在水泥地面上。影子在动。不是他在动。是风把灯吹得微微晃动。

他坐了下来。坐在地上。

房子里只有他一个人。蚊子三两只，在灯光周围转。江水的声音从门外传进来。门没关。

他坐了很久。不是在想事情。只是坐着。

天黑透以后，他锁了门，开着面包车回去。


\scene

第四天，老唐来了。他不是来看方远的。他是在附近办事，顺路弯过来。他的车比方远的面包车干净得多。黑色的轿车，轮胎没有泥。

他站在码头上，把周围看了一遍。

这就是你要的地方。

嗯。

老唐指了指铁皮房子。里面有床吗。

没有。

那你晚上睡哪。

方远指了指面包车的后座。

老唐没说话。他拿出烟，点了一根。这地方他年轻的时候来过。那时候还有船。船工在码头上卸货，汗衫贴在背上，露出一道一道的肌肉纹路。他站在同样的位置，看见的却是不一样的东西。

他说，你一个人不行。

方远说，我知道。

老唐说，你需要人。

嗯。

需要车。

嗯。

需要关系。

方远看着他。

老唐说，我不是在给你建议。我是告诉你，你缺什么。缺的东西你自己得去找。我一个都帮不上。

方远说，我知道。

老唐把烟抽完。烟蒂没捻灭，扔进江水里。兹的一声。

他看着那缕青烟从水面上飘起来。

说了一句：那个交通局的意见函。

嗯。

你是不是还没去找。

方远没说话。

老唐说，那个字，比水务局的字值钱。你要早点去。人在位置上，说不定哪天就换了。换了，你还得从头来。

方远点了点头。

老唐走了。车开出那段土路的时候，方远听到他车底盘刮到石头的声音。尖锐的。像指甲划过黑板。


\scene

交通局在城东。

方远以前来过。以规划局的身份来的，开过几次会，坐在会议室里，面前有茶和矿泉水。今天是以前同事的身份来的。

他打了马德胜的电话。马德胜以前参加过他的一个论证会。方远发过言，马德胜在下面做记录。会后两个人聊过几句。不是那种私下交情，是一种同行之间的识别。你一眼能看出这个人是不是在琢磨事。马德胜看的眼神和别人不一样。

电话通了。

马科长。我是方远。

哦。方工。好久不见。

我从规划局出来了。

电话那头沉默了一下。大概两秒。

现在做什么。

搞了个小货运。方远说，有个意见函想请你帮忙。

又是一阵沉默。然后是马德胜的声音，比刚才低了些。

你来我办公室一趟。

下午三点。

方远到的时候，马德胜的办公室门半开。他敲了一下。进来。那声音比电话里更沙哑。烟抽多了。

马德胜坐在桌子后面。办公桌上摊着三张报纸，一杯浓茶，一个烟灰缸。烟灰缸满了。

坐。他指了指对面的椅子。

方远坐下。椅子是折叠椅，坐上去比想象的软。

马德胜看着他，没说话。他在等方远开口。

方远把事情说了。渡口的位置，想做的业务，需要意见函。他没有说太多。他知道说得越少，对方越容易点头。说得太多，反而给了对方拒绝你的材料。

马德胜听完了。他端起茶杯喝了一口，放下。

有没有别的办法。方远问。

马德胜看了他一眼。不是审视，是一种重新打量的眼光。他见过很多来找他办事的人。有人求，有人暗示，有人带东西来。他都不喜欢。方远什么都没带。连请求都不是请求，是一个问题。

他用手指弹了一下烟灰。烟灰落在桌上，他用手指扫掉。

他拉开抽屉。抽屉里面很乱。印章、硬币、打火机、名片。他翻了半天，从最底下翻出一张名片，放在桌上。

这个。他说。

方远拿起来看了一眼。省厅的一个副处长。姓黄。

你说你搞的是小货运。

嗯。

马德胜说，小货运不需要这个级别的章。但如果你以后想做大，这个码头的定位要从临时装卸改成正式泊位。那就要省厅批。

方远把名片收进口袋。

马德胜又说，你先把区的章拿了。你这个位置不占航道，区里能批。章拿到了，先用着。上面这个，你以后再跑。

方远说，好。

马德胜站起来。他个子不高，站直的时候比方远矮半个头。但方远觉得他在办公室里比在外面高。

他走到门口。忽然停住了。

你在规划局的活儿，是不是干得不舒服。

方远说，挺舒服的。

那为什么走。

方远没有马上回答。他看到窗外的树。树叶子黄了。不是深黄，是那种介于青和黄之间的颜色。再过两天就全黄了。

他说，可能就是太舒服了。

马德胜看了他两秒。点了根烟。摆了摆手。

走吧。


\scene

区交通局的章，等了十天。

这十天里，方远干了很多事。打扫房子，给码头补裂缝。他买了水泥和沙子，在码头上和了三天。第一天和得太干，裂缝补上去又裂开了。第二天和得太稀，不凝固。第三天才刚好。多像一个工程。他想起以前看的地基报告。水灰比，骨料配比，养护时间。有些东西坐在办公室里翻图纸永远学不会。

他还买了一艘旧货船。花了六万。船身的漆掉了一半，露出下面绿色的一层。那层绿色的漆更老，可能是前船主刷的，也可能是前前船主刷的。发动机是大修过的，启动的时候冒黑烟，要预热五分钟。但能动。

卖船的是个五十多岁的男人，姓蔡。老船工，干不动了，儿子不接班。蔡师傅把船钥匙给他的时候，手颤了一下。不是病，是年纪。他看方远上船试了一圈，回来的时候说了一句话。

你会开。

会一点。

小时候上过我爹的船。

蔡师傅点点头。没再说什么。

船的名字叫盛丰三号。方远没改。

他把船开回了渡口。靠岸的时候没控制好速度，船帮撞了一下码头。声音很大。水泥上又多了一道新印子。他站在船头上看了看。拿了块石头在印子旁边画了个人脸。画完发现不像人脸，像一只猫。

他笑了。四下无人。


\scene

第十一天，章到了。

方远拿到了区交通局的意见函。函上只有一个段落，三十几个字。大意是经核实，此处不占用航道，同意备案。

他把这张纸小心地折好，放进外套里面的口袋。

然后他给鲁建国打了个电话。

电话响了很久。方远以为他不会接了。

最后一声的时候通了。

建国。

方远。

章拿到了。

电话那头安静了一小会儿。方远听见鲁建国身后有人在说话。像是在办公室。

鲁建国压低了声音。说，那个渡口。

嗯。

你真的搞起来了。

嗯。

又是一阵安静。

方远说，你上次说，你觉得我会做成。

嗯。

现在是了。

鲁建国没有说话。方远听到他那边有人催他开会。声音不大，但很清晰。鲁科长，会议开始了。

鲁建国说，我下了班去找你。

好。

方远挂了电话。

他坐在码头上。天开始下雨了。雨不大，落在江面上像针。一圈一圈，很快就不见了。落在脸上是凉的。他没有躲。

他在想，鲁建国会不会来。

不是在别的意义上想。是在字面意义上。他想看他会不会来。


\scene

鲁建国来的时候天已经黑了。

他的车在雨里很慢。雨刮器在前窗上来回摆动，车灯照出去只看到雨线的影子。他熄了火，没马上下车。坐在驾驶座上，看着雨中的码头。铁皮房子里亮着灯。一盏四十瓦的黄灯。方远站在门口。

他推开车门，跑了三步。衣服前襟湿了。

你疯了。鲁建国说。语气不是很重。像在陈述一个事实。

方远把一条干毛巾扔给他。毛巾上有沙子，干的已经结成小疙瘩。鲁建国擦了脖子。

两个人坐在门外。屋檐不够宽，雨飘进来，打在地上。鲁建国把腿缩了缩。

他没提老婆。方远也没问。

他看了一圈。铁皮房子亮了灯。码头上停着船。船上有一盏小灯，忽明忽暗。

他说，就你一个人。

嗯。

忙得过来吗。

还忙不过来。方远说，但我能忙过来。

鲁建国转过头来看他。方远在看雨。鲁建国也看雨。

雨落在江面上，看不清形状，只能听到声音。

鲁建国说，你缺人。

嗯。

缺一个能做事的。

嗯。

鲁建国说，我下个月能出来。

方远看着他。雨水从屋檐滴下来，断成一截一截的。

他说，你不是说你不会现在来吗。

鲁建国说，我说的时候是认真的。现在是现在了。

方远说，他们裁你了。

鲁建国没说话。他低下头。雨在地上积了一小滩水。他在滩水里看到自己脸的影子。很模糊。雨水打在脸上，他就碎了。

方远说，什么时候的事。

就这几天。鲁建国说，通知还没下来。但我已经知道了。

方远站起来。

他走进铁皮房子，从铁柜里拿出半瓶白酒。是之前蔡师傅给的。他说是用来暖船的，剩了半瓶。

他拿出来，拧开盖子，递给鲁建国。

鲁建国接过来，对着瓶口喝了一口。呛了。咳了两声。

好辣。

嗯。

鲁建国又喝了一口。这次顺了。

他说，你说你存了六年。

嗯。

我存了二十年。鲁建国看着手里的瓶子。二十年，一个通知。

方远没说话。

鲁建国把瓶子递回去。方远喝了一口，辣得皱了一下眉。

鲁建国笑了。笑的声音不大，从鼻子里出来的。

他说，你知道我最佩服你什么。

方远摇头。

鲁建国看着雨。

你走的那天，一个人在路灯下站着。手里提个纸袋。我说你接下来怎么办，你说走走看看。我当时觉得你是没想好。后来发现你是早就想好了。你只是不解释。

方远没接话。他又喝了一口酒。

鲁建国说，这个地方再过一年，水路货运的量能翻三倍。你那个时候就看到了。

方远说，水还在流。

鲁建国说，你又说这句话。

方远笑了笑。那个短的笑容。

他说，因为这句话是对的。

鲁建国把瓶子拿过去。最后一口。仰起头的时候，有雨落进他的嘴里，和酒一起吞下去了。

他站起来，把空瓶子放在门框旁边。它直直地立着。

他说，我今天不算正式来。等通知下了，我来找你。

方远说，好。

鲁建国走了。他的车在土路上颠了两下。尾灯划过芦苇丛，红了一下，又红了一下。然后拐弯了。

看不见了。

方远坐在门框下。

雨还在下。不急，也不慢。均匀的，像有人在数数。

他闭上眼睛。

他听见江水的声音。江水的声音是不变的。下雨也好，刮风也好，它的频率是一样的。

他想起了他父亲。父亲说，河不会骗人。

他想，下一句是什么。

还是没想起来。

他在雨声中继续坐着。

坐到灯灭了。柴油发电机自动停了。他没去修。

黑暗里，只有水声。

\clearpage

% ==== Chapter 3 ====
\chapter{旧厂}
\indent
化工厂在江城西边。

方远开车去的。面包车里的收音机坏了，只有一个频道能听。是个地方台，放老歌。他没有关掉。

路上经过老工业区。街道窄，两边是退了色的招牌。有一家饭馆，门开着，里面没开灯。招牌上写着三十年老店。字是红漆写的，被太阳晒得发粉了。

拐进一条土路，路边长着野草。草长到膝盖那么高，有一部分已经被人踩倒了。压出一条弯弯曲曲的小路。有车走过。

化工厂就在路尽头。

远远就能闻到味道。不是那种刺鼻的工业味，是另一种。像什么东西沤在太阳底下，发酵了，但又没有彻底腐掉。

厂门口没有人。铁门开着一半。门上的漆掉得厉害，露出下面黑色的铁皮。铁皮生了锈，锈是红色的，像干涸的血迹。

方远把车停在门外，走进去。

厂区很大。空地占了三分之二。地上有裂缝，缝里长着杂草。有一辆卡车停在角落里，轮胎瘪了，车身上落了厚厚一层灰。挡风玻璃裂了一条缝，用透明胶带粘着。

车间有一半开着门。里面没声音。机器停着，安静得像睡着了。

往里走。

办公室在最里面。一栋两层的小楼。楼前面有一棵梧桐树，叶子落了一半，没落完的那些也开始黄了。黄得不均匀，一半深，一半浅。

他在楼下喊了一声。

有人吗。

停了半分钟。二楼的窗户开了。

一个老头探出头。瘦，脸上皱纹很深，像刀刻的。他看了方远一眼。

你找谁。

方远说，我找赵厂长。

我就是。什么事。

想谈谈生意。

老赵看了他几秒。关上了窗户。

方远在楼下等着。没站多久，老赵下来了。穿着一件旧夹克，拉链拉到一半。手里端着一个搪瓷杯，杯子的白漆掉了一半，露出里面黑色的底。他走路不快，但每一步都很稳。

你是谁。老赵问。

方远。搞船运的。

船运。老赵重复了一遍这两个字。看了他一眼。

跟我上来。


\scene

办公室在三楼。

楼梯很窄，扶手上全是灰。每踩一步，灰尘就从台阶缝里飘起来。方远走得很慢，踩实了才往前。

办公室门开着。里面有一张桌子、两把椅子、一排旧柜子。墙上挂着营业执照，镜框很厚，玻璃反光。执照上的名字是江城化工厂，注册时间是一九八七年。还有一张奖状，是区里发的，日期模糊了。

桌上放着算盘。木头框子，紫檀色，珠子磨得发亮。

老赵把搪瓷杯放在桌上。坐吧。

方远坐了。

屋里只有两把椅子。方远坐了一把，另一把空着。老赵站着，从柜子里拿出一个暖水瓶。倒了杯水推过来。水是温的。

然后他问。

你多大。

三十二。

老赵点点头。像是在确认什么。

他坐下了。椅子发出吱嘎一声。椅腿在地上磨出一道浅印。

说吧，什么生意。

方远说，化工厂还有多少产量。

老赵看了他一眼。

不多。老赵说这话的时候没有表情，像在报天气。三成。有时候两成。

产品运到哪。

省内为主。有时候出省。

谁运。

老赵没回答。他端起搪瓷杯喝了一口水。盖上盖子。放下。

方远知道他不想说。每个厂都有自己的运输渠道，有的是固定的，有的是私人的，有的不好摆到台面上讲。

他说，我不是来打听你客户的。

老赵没说话。

方远说，我是来谈价格的。

老赵这才抬起头。你报个价。

方远报了一个数。

老赵把搪瓷杯放下。杯子磕在桌面上，发出一声闷响。

太低了。

这个价格是行的。方远说，我没压你。

老赵说，你那个价，连成本都包不住。

方远说，我的成本比别人的低。

老赵没接话。他站起来，走到窗边。窗外是厂区。空地、停了机器的车间、一辆落灰的卡车。

他说，这个厂一九八七年建的。那时候我二十三。

方远没说话。

那时候没人懂化工。老赵说，我跟一个浙江师傅学了三年，才把这个厂弄起来。中间死过两回。非典那年死了一回，零八年金融危机又死了一回。都缓过来了。

方远说，我知道。

老赵转过身看着他。你知道什么。

方远说，你还在做，说明还能做。

老赵盯着他看了很久。那种眼神像是在掂量一个人的重量。他不是在掂量钱，他是在掂量这个人。看他值不值得聊下去。

坐吧。老赵说。我再倒杯水。


\scene

水喝了一壶。又一壶。

方远没有再报价。老赵也没有再问。

他们聊了一些别的。化工厂的事，老客户的事，运费结算周期的事。老赵说话很慢，每个字都像是从很深的地方捞上来的。方远听着，偶尔应一声，不多话。

太阳从窗户的一边移到另一边。光线越来越斜，落在墙上，把那面墙切成两半。一半亮，一半暗。

快三点的时候，老赵站起来。他走到墙角，打开一个柜子，拿出一个纸箱。纸箱里是旧账本，封面上写着年份和月份。一本一本，码得整整齐齐。

他把最上面一本拿出来，翻到某一页，给方远看。

这是上个月的出货记录。

方远看了一眼。数字不多。出货量是三年前的三分之一。

老赵说，不是没订单。是运不出去。

方远抬起头。

老赵说，现在的运输公司要现结。我这边账期压着，一压压三个月。资金周转不过来。订单接了不敢做，做了运不出去。

方远没说话。

老赵合上账本。把它放回纸箱里。

你报的那个价。老赵说，能做。

方远看着他。

但你要接受账期。老赵说，三个月。

方远说，不行。

老赵把纸箱放回柜子里。关上门。转过身。

你知道这三个月压的是多少钱吗。

知道。

你做得了主吗。

方远说，做得了。

老赵看着他。这次看得更长了。

然后他坐下了。

再喝杯水。

第三壶水是四点的时候倒的。水温不对，太烫了。老赵喝了一口就放下了，说水太烫。方远也喝了一口。太烫了。但他还是喝完了。


\scene

他们聊到五点半。

太阳下山了。屋里暗下来。方远站起来去开灯。开关在门边，他按了一下，灯没亮。

停电了。老赵说。

方远在黑暗里站了一会儿。他掏出手机，打开手电筒。

光打出来，照在老赵脸上。皱纹很深，颧骨很高。

老赵没躲。他说，这个厂经常停电。夏天最严重。有一次停了三天。

方远把手电筒放下。手机的光弱，照不远。屋里只有一小块是亮的。

他往窗边走。站到窗边。窗外也暗了。厂区的灯亮了几盏，不多。卡车还停在那个角落，轮廓模糊了。

他突然看到墙角有一台落地扇。黑色的扇罩，铁支架，底座很沉。那种老式的风扇，现在很少见了。

坏了？方远问。

老赵说，坏了两周了。

方远走过去。蹲下来看了看。扇罩卸了一半，底座里露出里面的电容。电容烧了。旁边的螺丝刀还在，工具箱打开着，里面有老虎钳、起子、扳手，一样一样码着。

你修的？

老赵没说话。他站在原地，看着方远蹲在那台风扇前面。

方远伸手去摸电容。烧过的那种电容，顶部会鼓起来。这台的顶部没有鼓，但有一圈焦痕。烧得不严重，但不能再用了。

你有备用的吗。

柜子里有一个。老赵说，但不知道能不能用。

方远站起来。走到柜子边，打开。一个纸盒，纸盒里有两个电容。他拿出来看了看。一个是好的，参数对得上。

他拿了一小截电线。两根铜丝。他把电容接上，试了一下。通了。风扇转了一下，抖了一下，又停了。

方远重新接了一下。再转。这次风扇转起来了。扇叶转得很慢，像老人下楼梯。但它在转。

方远站起来。

退后一步。

老赵走过来，看着转动的扇叶。扇叶带起一点风。很微弱，但能感觉到。

老赵站在那里，看着。

屋里很安静。只有风扇转动的吱吱声。很细，像虫子在叫。

站了很久。

老赵开口了。

方远。

嗯。

三个月账期，我给你减一个月。

方远看着他。

两个月。老赵说，这是我的底线。

方远没有马上答应。他站在那里，看着那台转动的风扇。扇叶转得很慢，每一圈都能听见声音。

他忽然问了一句。

这台风扇，多少年了。

老赵想了想。说，三十二年。

方远说，比这个厂还早。

老赵说，是我父亲那辈人置办的。那时候厂还是我父亲管。我接手的时候，电风扇还在。我接手以后，换过一次电机。其他没动过。

方远点了点头。

他看着老赵。

两个月。行。

老赵伸出手。两只手握在一起。老赵的手很粗糙，指节粗大，像干了一辈子体力活的人才有的手。

握完，老赵转身把风扇关了。扇叶慢慢停下来。屋里又安静了。

你吃点东西再走。老赵说。

不用了。

吃点。老赵说，我这儿有面。


\scene

面是老赵下的。

锅是那种老式铁锅，灶台是水泥砌的。切菜用的是菜刀，刀背磨得发亮。青菜是院子里种的，葱是院子里长的。

面端上来的时候，窗外已经完全黑了。

老赵坐在对面。桌上还有一瓶酒。不是好酒，普通的白酒，玻璃瓶。他倒了两个半杯。

方远端起来。

老赵说，我不问你的事。

方远说，嗯。

但你这个人，我看懂了。

方远看着他。

老赵把酒喝了。喝得很慢。放下杯子的时候，筷子夹了一块腐乳。

这个厂，我儿子不想接。老赵说，他在深圳做IT，一年挣的比这个厂十年都多。他不回来看我，我也不指望他回来。

方远没说话。

老赵说，我六十三了。干不动了。但我不想卖。卖了就什么都没有了。

方远说，我懂。

老赵看了他一眼。

你不会把它卖给别人的。

不是问句。是陈述句。

方远说，不会。

老赵点点头。又喝了一口酒。

窗外的黑暗里，有灯一闪一闪的。不知道是什么。

吃完面，方远站起来。

老赵没送他下楼。只站在门口，看着他走到楼梯口。

方远回过头。老赵站在三楼办公室门口，身影被走廊的灯光拉得很长。矮胖，像一截树桩。

明天把合同送过来。老赵说。

好。

方远走下楼梯。楼梯间很黑，他打开手机的手电筒，一点一点往下照。扶手上的灰尘比来的时候少了。

走出厂门的时候，他看了一眼门上的牌子。江城化工厂。五个字，红色，漆落了一半。

他上车，发动引擎。

收音机还是那个台。还是那首老歌。他调大了声音，歌听不清了，只剩下调子。

他把车开动。

后视镜里，厂区的灯越来越远。最后缩成一点，然后不见了。


\scene

回到渡口是九点。

鲁建国在铁皮房子里等他。桌上放着两盒快餐，还温着。

你又不吃晚饭。鲁建国说。

方远坐下，把快餐打开。是盖浇饭，土豆烧肉。肉有点老。

他吃了一口。

鲁建国坐在旁边，也在吃。两个人没说话。门开着，江风灌进来，把灯光吹得晃了晃。

吃到一半，鲁建国问了一句。

谈下来了？

谈下来了。

两个月账期。

鲁建国停了一下筷子。

你吃亏了。

方远没抬头。继续吃。

他说，老赵这个人，值这个价。

鲁建国看着他。没说话。

方远把最后一口饭扒进嘴里。站起来，把餐盒扔进门口的垃圾桶。

他站在门口，看江。

鲁建国也站起来。走到他旁边。

老赵的厂，还有多少产能。

三成。

够吗。

不够。方远说，但够了。

鲁建国没听懂。但他没有再问。

他进了屋子。

方远站在门口很久。江水的声音从黑暗里传过来。和白天一样。和昨天一样。和明天一样。

他想起了老赵屋里那台电风扇。转了三十多年，还在转。

他想起了老赵的手。粗糙，像老树皮。

他想起了老赵说的那句话。

你这个人，我看懂了。

他笑了笑。那个短的笑。

然后他也进了屋子。

门关了。

柴油发电机在响。嗡嗡的，像另一种声音的呼噜。屋子里亮了。灯是四十瓦的黄灯，光落在水泥墙上，像一层旧纸。

他躺下来。床是一张行军床，钢丝的，躺上去有声音。他听着那个声音，然后闭上眼睛。

江水在门外流着。

他睡着了。

\clearpage

% ==== Chapter 4 ====
\chapter{大雨}
\indent
雨是三月开始的。不是那种突然下来的暴雨，是慢慢的、一层一层的。今天比昨天多一点，明天比今天再多一点。到你发现的时候，已经下了好几天了。

那是方远做货运的第五个月。账上多了三家客户。除了老赵的化工厂，还有对岸一家水泥厂和一家纸箱厂。水泥厂的货最重，一船下来船舷压到吃水线。纸箱厂的货最轻，轻得船在水上晃，像一段木头。

三个客户，三条航线。渡口从这里出发，沿着江往下游走。十五公里到水泥厂，三十公里到纸箱厂。最远的是老赵，在城西，回去的时候要逆水走一个小时。

方远每天五点起床。柴油发电机打着，烧一壶水，泡一杯茶。然后上船。发动机预热的时候，他坐在船头借着天光看水位。水在涨。

三月的雨把整个江面抬高了一米。平时露出水面的礁石大部分被淹了。看不清水下的东西，船有时候会擦到暗礁，刮出一个深沉的闷响。他不怕撞击。他怕螺旋桨被水草缠住。那种缠是悄无声息的。螺旋桨还在转，但船不动。你回头看，拖着一长条绿色的东西。

他已经停过两次。都是小雨里，趴在船尾用刀割水草。水草比绳子还韧，割断了就沉下去，下一段缠上来。


\scene

雨一直下。一天接一天。方远的手上全是裂口。不是割伤的，是被水泡出来的。皮肤泡涨了，泛白，然后裂开。他不看伤口，该干什么干什么。疼是有的。但淋在雨里，那种疼反而让他清醒。冷也是一个道理。

有一天傍晚，他送完水泥厂的货回来。逆水。船走得很慢，发动机的声音比平时沉。江水是浑的，浑黄色，卷着泥和断枝往下冲。他看着江面上漂下来的东西。一只塑料拖鞋，一个破渔网，一截泡涨的树干。拖鞋漂过去的时候，他不知道为什么多看了两眼。是蓝色的，左脚。

雨落进水里，没有声音。不是没有声音，是被江水的声音盖住了。

船到渡口的时候天已经黑了。他拴了缆绳，走进铁皮房子。头发和衣服都湿透了，往下滴水。他脱了外套，拧了一下，地上就是一滩水。他把茶缸放到电炉上烧。等的时候浑身发抖。不是怕，是身体的反应。他管不住。

水烧开了，他的牙齿还在抖。

他喝了口热茶，牙齿才慢慢停下来。


\scene

鲁建国的电话是那天晚上打来的。

方远刚换了干衣服，坐在床上。柴油发电机的声音盖过了手机铃声。他听到的时候手机已经响了不知道多少声了。嗡嗡的，像一只苍蝇。

他接起来。

建国。

那边是沉默。方远听出来不对。鲁建国的沉默有两种。一种是在想怎么说，一种是不想说。这次的沉默是后一种。

过了好一会儿，鲁建国说，下了。

方远说，什么时候。

今天。

方远没有说话。他知道这一天会来。当时鲁建国告诉他通知还没下来但自己已经知道了，他就知道这一天会来。但他没想到会在雨夜。在这个铁皮房子里，伴着江水的声音，听到这句话。

我不走。鲁建国说。他那边也有雨。方便听到雨打在什么东西上的声音。可能是窗台，可能是空调外机。

方远说，你在哪。

家外面。

鲁建国的声音很低，像是在躲着什么人。方远知道他在躲谁。他老婆。他老婆从一开始就不相信他。不是不相信方远，是不相信任何需要把存了二十年的钱押上去的事。

方远说，来帮我。

他说的时候声音很平静。没有加重，也没有放轻。就是陈述的一句话，像说我今天多吃了一碗饭。

电话那头的沉默比前面都长。

方远听到雨声。他这边也有，电话那头也有。两种雨声混在一起，分不清哪个更急。

鲁建国说，我老婆骂了我一晚上。

方远说，嗯。

她说我年纪大了，出去能干什么。

嗯。

她说你会坑我。

方远没有说话。他把手机换到另一只手。右手已经被他的体温焐热了。他换到冰冷的左手。左手不抖了。

鲁建国说，我跟她说，不是方远坑我，是我自己签的字。

方远说，你说得对。

鲁建国说，但她还是不同意。

又是一阵沉默。

然后鲁建国说，给我两天。

方远说，好。

他挂了电话。手机搁在旁边的椅子上。屏幕暗了。铁皮房子里只剩下雨声和发电机声。

他站起来。走到门口。拉开了门。

雨很大。是从另外一个方向来的。没有打在他脸上，但他能看清楚每一滴雨落在地上的瞬间。水泥地面上积了一小滩水。雨滴打在水面上，起了无数个水泡，很快就破了。

他看着这些水泡。看了一会儿。然后关上门。


\scene

第二天雨还是没停。

方远照常出船。水泥厂的货。在江心的时候他看见对岸有人在跑。跑得很快，用外套蒙着头。看不清是谁。可能是船工，也可能是附近的居民。他看了一会儿，那人跑进一条巷子里，不见了。

船到水泥厂。码头上的工人戴着斗笠，雨打在塑料布上噼里啪啦地响。方远帮着卸货。肩膀湿了，热气从衣领里冒出来。工人的手和他的一样，全是裂口。

卸完货，工人递给他一支烟。他接过来，工人帮他点。雨里打火，试了三次才着。他吸了一口，咳了两声。他已经很久没抽烟了。上次抽还是鲁建国留给他的那包。那根被他扔进江里了。

回去的时候，他特意绕了一段路。往上走了五公里。那里有个水文站。是一间很小的房子，蹲在江边，像一只鹭鸶。

他把船靠过去。一个老头坐在门口，看着雨。

水位怎么样了。方远问。

老头看了他一眼，指了指门口的标尺。方远上了岸，走到标尺前面。

水已经涨了两米。标尺的最上面一段还露在外面，但按这个速度，再下两三天，最上面也会淹掉。

还涨吗。方远问。

老头说，报的是还要下。上面那个水库已经泄洪了。下来的水还没到这里。最晚明天夜里。

方远在水文站门口站了一会儿。雨从屋檐上流下来，在他脚边淌成一条小河。

他说，谢了。

老头没应。他已经在下一支棋了。棋盘摆在膝盖上，棋子被雨打湿了一小块。


\scene

那天晚上他没睡。有一种感觉，说不清是什么。不是焦虑，也不是预感。就是知道有事。

他坐在床上，靠着铁皮墙。墙上有一小片地方在渗水，摸上去是湿的。他把枕头移了一点，不让它碰到水。

凌晨三点，雨更大了。屋顶的铁皮开始响。不是那种均匀的节奏，是一阵一阵的，像有人在屋顶上泼水。他听见远处的雷声，很远，沉沉的。

四点半的时候电话响了。

他接起来。不用看也知道是谁。

鲁建国的声音。他说，我在路上了。

方远坐直了一点。

在哪条路。

土路。雨太大，看不清。

你怎么来的。

开车。只带了两个箱子。她没醒。

方远说，开慢点。

鲁建国说，我开到进来土路。前面有一段被水淹了。我下车看了看，能过。

别冒险。

鲁建国说，已经过了。

五点半，雨小了一些。方远在门口等着。手里握着一把手电筒。

灯光从远处照过来。先是一个白点，然后白点变大，然后是整个车头。车在土路上走得很慢。轮胎陷下去又弹起来，溅起泥水。

车停在码头旁边。车灯照在方远身上，他一动不动。车灯关了，黑暗重新合拢。

鲁建国从车上下来。

他穿了一件旧雨衣。雨衣太小了，袖子短了一截。手腕露在外面，全是雨水。他没撑伞。

方远看着他。

两个人对视了几秒。

方远说，进去。

鲁建国进了铁皮房子。方远把门关上。外面的雨声小了。

鲁建国环顾了一圈。一张行军床，一张桌子，一把椅子，一个铁柜子，一个电炉，一个茶缸。墙上渗着水。

他说，你就住这。

嗯。

我两个箱子，放哪。

方远看了看房间。比你家里小一点。

鲁建国说，小很多。

方远笑了笑。那个短的笑。

鲁建国也笑了一下。他笑起来很勉强，像嘴角被什么东西扯了一下。把雨衣脱了，叠成一个不规则的方块，放在桌脚旁边。他蹲下来，把雨衣的一角压平了。

他站起来。

鲁建国说，我这个人没有什么本事。一辈子就在那个办公室里看新闻、填表、开会。被人裁了，才觉得以前那些年白过了。

方远没有说话。

鲁建国说，但我会做事。你让我做什么我就做什么。船工也行，搬货也行，看门也行。

方远说，先把头发擦干。

鲁建国接过毛巾。毛巾是半干的。他擦了擦头。毛巾上有很熟悉的味道。洗衣粉和江水混在一起的味道。他在规划局的洗手间闻了六年。现在在这里又闻到了。

鲁建国把毛巾还给他。

他说，你现在一共有几条船。

一条。

几个客户。

三个。

鲁建国沉默了一秒。说，这么大的雨，你一条船够吗。

方远说，不够。

那你还说行。

嗯。方远说，因为只要水还在流，船就能走。

鲁建国看着他。说，你每次说这句话，我都不知道你是什么意思。

方远说，我也不知道。

天还没有亮透。雨小了很多。能听见雨打在铁皮上的声音变小了，从泼水变成了滴水。一滴，一滴，间隔越来越长。

他开了门。河面上有一层薄薄的雾气。看不清远处。但水声没变。不管雨大雨小，水声是一样的。

他说，今天不出船。

鲁建国说，为什么。

方远说，水位还在涨。上面泄洪了。最晚夜里到。

鲁建国哦了一声。沉默了一会儿。

他说，现在做什么。

方远指着门外的码头。

修码头。


\scene

雨停了半天，又下了半天。

他们穿着雨衣在码头上干了六个小时。把所有的裂缝补了一遍，加固了缆桩。鲁建国的手比方远嫩。干到两个小时的时候，他的手就开始起泡了。他没吭声。继续搬石块。泡破了，血从手套里渗出来，把白色的手套染成暗红色。方远看到了，没让他停。只是拿了自己的手套，把他的换了。

午饭是泡面。在铁皮房子里吃的。电炉上烧水，两碗面泡开了。鲁建国吃了一口，说，这面和你六年前请我吃的那碗一个牌子。

方远想了想，说，是吗。

是的。那天你刚来规划局。我请你吃的。也是这个牌子的泡面。

方远说，我以为是楼下卖的盒饭。

不是。是泡面。红烧牛肉味的。跟现在这碗一样。

方远看着手里的泡面杯子。红色的包装。上面印着几块酱色的牛肉，油亮亮的。他没再说什么。只是低头吃面。

汤喝完了。他把杯子放在桌上。

两个人又出去了。下午雨又大了一些。他们修到天黑，实在看不见了，才收工。

码头还是旧的。裂缝也还在。但不会再继续裂了。方远看着鲁建国，他满手是血，坐在门框上不说话。

方远说，谢谢。

鲁建国没回头。他说，不客气。

停了很久。

他说，我老婆说你那句是最对的。

哪句。

她说你说走走看看的时候其实早就想好了。方远没有说话。风从门缝里灌进来，吹得电灯晃了一下。

鲁建国说，我现在明白了。你不是走走看看。你是在等。

等什么。

等水位上来。

方远转过头来看着他。他蹲在地上，手上包着纱布。纱布上有血渗出来，在黄色的灯光下变成褐色。他没有回答鲁建国的问题。

把铁门关好。他说。

供电断了一阵。不知道是哪里出了问题。可能是发电机本身的毛病，可能是油不够了。鲁建国去检查了。两人用手机的光照着，修了二十分钟。配电器的接口松了，拧一下就好了。灯重新亮起来的时候，两个人互相看了一眼。方远看到他头上的汗。他没有擦。

那天晚上鲁建国睡在屋里。东西都还没弄好。铁皮房子里只有一张行军床和一个睡袋。窗户已经破了，以前用工地上捡的木板钉死了，但光线还是从缝隙里漏进来。外面的雨声很远，像是隔了很多层东西。

睡袋很薄。鲁建国裹在里面，像一只虫子。

方远，我睡不着。

嗯。

你在想什么。

方远说，我在想船。雨一停，江水的流速会加快。起锚的时间要提前。不然水会把船往下游推，推得太远，收不回来。

他说这话的时候，语气很平常。

鲁建国听着，翻了个身。行军床发出的声音很大。

他也想着那条船。花六万块买的，发动机预热要五分钟。漆掉了，高油。但它在江上走。

鲁建国闭上眼睛。

雨还在下。

但他睡着了。

\clearpage

% ==== Chapter 5 ====
\chapter{许可}
\indent
雨停了三天。江水开始退。

方远在码头上看着水位线。退了一米二。还不够。但够让船走起来。

鲁建国已经在渡口住了两个星期。他每天早上五点跟着方远起来。方远开船，他在码头上等着。船回来了，他卸货。货卸完了，他搬进铁皮房子里码好。铁皮房子的一半已经被货物占了。纸箱堆到天花板，化工原料用防水布裹着，扎得紧紧的。空气里有一股淡淡的化学品的味道，像氨水但不刺鼻。

两个人说话不多。方远开船的时候，鲁建国在岸上做事。方远回来的时候，鲁建国把东西搬好。他们之间的默契不是用话堆出来的。是一个人做了什么，另一个人看见了，下次就不用说。

方远在心里算了一笔账。三个客户的量加起来，一个月大概能跑十五趟。每趟的利润不多。扣掉柴油、维修、生活费，剩下来的刚好够维持。不够扩展。

要扩展，就要接更多的货。接更多的货，就要更大的船。更大的船，就要正式的泊位批文。

临时装卸的许可是区里给的。正式泊位要省里批。

他想起了马德胜给的那张名片。


\scene

省交通厅在省城。开车要三个小时。

方远选了个晴天上路。面包车的收音机还是那个频道，但出了江城的地界就收不到了。他把收音机关了。路上只有发动机的声音和轮胎碾过路面的声音。

省城比江城大三倍。路上车多，堵了两次。他找到一个停车位，停了车，步行去了交通厅。

交通厅的大楼很新。玻璃幕墙反着光。门口有保安。保安看了他一眼，问找谁。他说了名字。保安查了一下，让他登记。他写上方远两个字，在来访事由一栏写了咨询。

电梯到八楼。走廊很安静。地上铺着灰色的地毯，走上去没有声音。他在803门口站了一下，敲门。

进来。

推开门。办公室不大，但很亮。窗户对着街道，阳光照进来。桌上很整齐。一摞文件、一台电脑、一只陶瓷杯。

黄副处长坐在桌后面。四十出头，头发梳得整齐。穿着白衬衫，没有打领带。袖子卷到小臂中间。

方远？

是我。

坐。黄副处长指了指对面的椅子。

方远坐下。椅子很软，比老赵那把折叠椅舒服多了。他不太习惯。脊背没有完全靠上去。

马德胜跟我提过你。黄副处长说，你是规划局出来的。

出来了。

自己做的？

做货运。在江城。

黄副处长点了点头。他端起杯子喝了一口水。放下杯子的时候，手指在杯壁上停了一秒。

方远注意到他的手。手指很长，指甲剪得很干净。不像做具体业务的人的手。

你说你想批一个正式泊位。黄副处长说。

嗯。

在江城哪个位置。

方远说了那个渡口的位置。黄副处长没有去翻材料。他在听。听完了，沉默了一下。

那个位置我知道。他说，以前是渡口。后来桥修好了，渡口就废了。

嗯。

你要在那里建正式泊位。

方远说，是。

黄副处长靠在椅背上。他看着方远。那种看法很特别，不是审视，也不是打量。是一种专业人士在评估另一个专业人士的感觉。

他说，这个事情有难度。

方远没有说话。他知道接下来黄副处长会说为什么有难度。他不需要问。他只需要等。

果然。黄副处长说了几个原因。第一，那个位置不在省里的港口规划范围内。第二，上下游都有已有的码头，间距不够。第三，河道在那一段有弯度，从航道上讲不适合大型船舶停靠。

方远听完了。

他只问了一句：有没有别的办法。

黄副处长看了他一眼。

这一眼比之前所有的看都重。方远懂这一眼的意思。这一眼是在说，你这个人到底是真的想做事，还是只是来碰碰运气。

方远迎着他的目光。没有躲。

黄副处长站起来，走到窗边。他背对着方远，看着窗外的街道。街道上有车经过，声音隔了玻璃听不到。

过了一会儿，他转过身。

有一个途径。他说，不属于正式审批流程。属于区域协调的范畴。如果江城和邻市的交通局联合出具一份区域水运规划调整的建议函，省厅可以据此做补充规划。补充规划里可以纳入新的泊位。

方远说，邻市？

黄副处长说，下游那个市。他们也在做水运振兴。如果你们的泊位能服务到他们的货物，他们有理由参与。

方远想了想。

他说，我认识邻市的人吗。

黄副处长回到桌前。他拉开抽屉，翻了一下，拿出一张纸。纸上有一个名字和一个电话。

这个人叫周学文。邻市交通局的。他管水运规划。你去找他。

方远接过来。看了一眼。记住了。然后把纸折好放进衬衫口袋。

黄副处长说，我不是在帮你。我是在告诉你路怎么走。路你自己走。

方远点了点头。

他站起来。谢了。

黄副处长摆了摆手。方远走到门口的时候，他又说了一句。

方远。

嗯。

你的衬衫有褶子。

方远低头看了一眼。

黄副处长说，下次来穿件熨过的。


\scene

去邻市的路上，方远在服务区停了一下。

他买了一个面包和一瓶水。坐在面包车里吃。面包是袋装的，口感像在嚼纸板。水是矿泉水，没有味道。

他看了看自己衬衫上的褶子。确实有。出门前用手抹了两下，没用。他不会熨衬衫。他也不知道去哪熨衬衫。

他把面包吃完。把包装袋塞进车门上的储物格里。发动车。

邻市比江城小。但沿江的那一段比江城热闹。有新建的码头，码头上停着两艘大型货船。船身上写着邻市港运的字样。吊臂在动。

周学文的办公室在交通局三楼。交通局是一栋旧楼，外墙上爬着爬山虎。十一月的爬山虎是红色的，像一面旗帜。

方远到了。敲门。

进来。

周学文比黄副处长年轻。三十多岁。戴着眼镜。说话快。

方远。

是我。

马德胜让你来的？不对，黄处长让你来的。

嗯。

方远把来意说了。周学文听的时候没有打断。他一直在转手里的一支笔。笔是黑色的，转得很快，没有掉过。

方远说完了。

周学文把笔放下来。

你说你的泊位在江城西边。那个位置，对吧。

对。

他说，我们的货走陆路到江城，再转运。如果你们有正式泊位，我们的货可以直接走水路到你们那里，再分拨。

方远说，省六十公里。

周学文点了点头。

他站起来。走到墙上的一张地图前面。地图很大，覆盖了半面墙。上面标注了所有的码头、航道、桥梁。

他指了指江城西边的那个位置。

这里。他说，如果我们出一份联合建议函，省里批了补充规划。你们拿到正式泊位，我们的货就可以走水路过来。

方远说，对。

周学文转过身看着他。

你要什么。

方远说，联合建议函。

我为什么帮你。

方远没有说因为这对你们也有好处。他没有解释。他只是看着周学文。

周学文也看着他。

过了几秒。周学文说，你们一个月的运量多少。

现在不多。方远说，但我可以保证你们走水路的成本比陆路低三成。

三成。

三成。

周学文想了一下。他走回桌前，坐下来。打开电脑，看了一个什么东西。方远没有凑过去。

周学文说，我需要看一下我们的运输数据。你等一下。

方远等着。窗外有一只鸽子落在窗台上，歪着头看了他一眼，又飞走了。

十分钟后，周学文合上了电脑。

他说，三成可以做到。但我有一个条件。

什么条件。

你的泊位要留一个专用泊位给我们。

方远说，可以。

周学文说，那我来写建议函。你回去等消息。

方远站起来。

周学文也站起来。他伸出手。方远握了一下。周学文的手不粗糙，但也不滑。是那种经常写字的人的手，食指侧面有一个茧子。

出去的时候，方远在走廊里站了一下。他看了一眼窗外的爬山虎。红色已经深了一些，快要变紫了。


\scene

回到渡口的时候天已经黑了。

鲁建国在码头上等他。船不在。方远出门的时候交代他把今天的货运完。他已经运完了。

卸了？

卸了。

方远点了点头。走进铁皮房子。桌上放着两碗泡面。他坐下来，拆开泡面，倒上热水。

鲁建国坐在对面。也在吃。

方远吃了几口。说，邻市交通局的人答应了。

什么答应了。

联合建议函。方远说，如果省里批了补充规划，我们就能拿到正式泊位。

鲁建国停下筷子。

正式泊位意味着什么。

意味着可以停靠五百吨以上的船。方远说，意味着可以接大单。意味着不只是老赵那点化工厂的货。

鲁建国想了一会儿。

他说，你哪来的这些关系。

方远说，以前在规划局的时候积累的。

鲁建国说，我以为你那些关系都没用了。

方远说，关系不是靠用的。是靠不用的。

鲁建国想了一会儿。他说，我没见过你私下跟谁吃饭。没见你送礼，也没见你请客。你是怎么让人帮你的。

方远把泡面的汤喝完。杯子放下来，烫手。他吹了一下手。

他说，你在办公室看新闻的时候，我在看图纸。你开会的时候打瞌睡，我在听每个人说什么。聊的不是工作，是他家孩子在哪上学，他老丈人住哪家医院。然后你需要一个东西的时候，你记得他不是因为他能帮你。他帮你不是因为你求他了。是你还记得他家孩子的事。

鲁建国看着方远。看了很久。

他说，你什么时候学会的。

方远想了想。在规划局学会的。他说。六年。每天看人。看图纸。看水。三个人看久了，有什么区别。图纸看久了，每条虚线什么意思。水看久了，什么时候涨，什么时候退。人看久了，谁说了算，谁该找谁，不用问。

方远把泡面吃完了。汤也喝了。

他说，你明天去一趟城里的干洗店。

干什么。

方远看了他一眼。把他衬衫上的褶子拉了拉。

熨几件衬衫。他说，下次去省城，不能穿皱的。

鲁建国看了看自己的衬衫。也看了看方远的。

他说，你自己的呢。

也熨。

鲁建国笑了。这次是真的笑。声音不大，但是从肚子里出来的。他把泡面吃完。把杯子叠到方远的杯子上。两个杯子叠在一起，站了起来。

我去洗碗。他说。

方远看着他走进码头边的小水房。水房没有屋顶，是一个砖砌的台子。上面有水龙头。鲁建国拧开水龙头。水是冰凉的，冲在杯子上，他缩了一下手。然后继续洗。洗完，把杯子倒扣在砖台上。两个杯子并排。

方远看着那两个杯子。风一吹，杯口的水往下流。

他说，好。明天去。

方远站起来，走到门口。

江面上有月光。不是满月。弯的。光落在水上，像一条银色的鱼。

他站在那里看了一会儿。

水在流。

月光跟着水流。流到看不见的地方，就没了。但水还在流。月光没了，水还有。

\clearpage

% ==== Chapter 6 ====
\chapter{过客}
\indent
陈敏的面馆开了三年。

说是面馆，其实就是江边一间铁皮棚子。棚子不大，摆得下四张桌子。桌子是折叠的，桌面上的塑料贴皮已经翘了边。天气好的时候门全打开，江风吹进来，比空调管用。天气不好的时候关上门，屋里全是面汤的蒸汽。

菜单只有四种面。素面、肉丝面、鸡蛋面、大排面。不写价钱。老客都知道。新客也不问钱，先吃，吃完再说。从没跑过单。

陈敏四十二岁。头发扎在脑后，用一个塑料夹子夹住。脸上的皮肤被江风吹得很干。不好看。但也不难看。是那种你看完就会忘掉的脸。

她一个人做所有的事。煮面、端面、收钱、洗碗。不忙的时候站着看江。忙的时候只有手在动。手很快。一把面丢进锅里，筷子搅两下，捞出来，浇汤，放浇头。一碗面从下锅到上桌，不超过三分钟。


\scene

那天的第一个客人是钓鱼的。

老宋。是这片的固定人物。每天早上六点准时出现。一个人，一根竿，一个小马扎。坐在江边那块凸出来的石头上，一坐就是半天。

要什么。陈敏问。

素面。

加不加辣。

不加。

老宋坐下来。陈敏把面端上来。他拿筷子的时候手有点抖。不是冷的，是老了。他自己也知道。筷子在汤里搅了两下，没夹起来。再搅两下，夹起来了。送到嘴里。嚼得很慢。

陈敏看着他。看了三年了。每次来都是素面，不加辣。话很少。问一句答一句。不问就不说。

今天钓了几条。陈敏问。

三条。

大不大。

不大。

老宋吃完了。把碗推到一边。从口袋里摸出一张钱，压在碗下面。站起来，拿着竿，走了。他走去那块凸出来的石头。坐下去。甩竿。线落在水面上，弹了一下，慢慢沉下去。

陈敏收了碗和钱。

老宋用的钱永远是旧的。软塌塌的，像布一样。她不知道他以前是做什么的。只听人说过他以前在江对面那个镇上住，后来桥修好了，他反而搬来了这边。问他为什么，他说那边太吵。

陈敏站在门口洗了碗。抬头看了一眼江面。老宋还在那里坐着。竿没动。水在流。阳光照在水面上，碎成无数个小光点。那些光点在浪尖上跳。跳得她眼睛有点疼。


\scene

上午来了几个学生。

四个。背着双肩包，脖子上挂着相机。看着像是美院的学生。头发都长。穿的衣服不修边幅。其中一个的鞋带没有系，拖在地上，踩脏了。

吃面吗。陈敏问。

吃。一个女孩说。她戴着圆形的眼镜，镜片很厚。

四个人点了四碗肉丝面。加了辣。吃着吃着开始聊天。聊的不是学校里的事。聊的是江水的颜色。说今天的江水偏灰，像水墨的淡墨。昨天是偏黄的，像熟宣的底色。

陈敏没听懂。她站在灶台后面听他们说话。他们说了很多词。什么肌理，什么质感，什么色调。都是陈敏不认识的字。

女孩吃到一半，端着碗走到江边。蹲下来。把碗放在膝盖上。一只手拿筷子，另一只手举着相机对着江面。

她拍了几张。然后把相机放下，低头吃面。筷子夹起来的面条往下滴汤。滴在她的裤子上，她没在意。

他们是来看江的。陈敏知道。

江边这几年来了很多这种人。拍照的。画画的。写东西的。他们觉得这条江美。那种美陈敏也看得见，但她形容不出来。她只觉得这江每天都一模一样。早上是这样，晚上也是这样。今天是这样，明天还是这样。

也许这就是他们觉得美的地方。一样的东西，你看久了就习惯了。他们看第一次，觉得新鲜。


\scene

下午来了一个穿衬衫的人。

衬衫是白色的，领子很挺，袖口的扣子扣得很紧。脚上穿着皮鞋。皮鞋擦得很亮，在土路上走了几步，就蒙了一层灰。

他走进面馆。环顾了一圈。不像来吃面的，像来考察的。

后面还跟着一个人。穿着夹克，拿着一个文件夹。夹克男一直在说话。说这附近的水路交通，说货运成本，说下游那个码头一年能赚多少钱。衬衫男听着。没有回应。只是看。

陈敏看出来了。是投资人。她见过这种人。前两年也有过几个。来看沿江的地块，谈水运和仓储。谈了几个星期，然后就不来了。不知道是没谈成，还是谈了就不需要再来。

两个人找了个位置坐下来。不是面馆的位置。是江边的一个石墩子。夹克男从文件夹里抽出一张纸，指给衬衫男看。衬衫男扫了一眼，又抬起头看江。他看的不是水。是对岸的开发区。那些刚封顶的楼房。那些塔吊。

陈敏走过去。

吃面吗。

两个人都愣了一下。好像是才意识到这是一个面馆。夹克男看了看衬衫男。衬衫男点了点头。

素面还是肉丝面。

都行。夹克男说。

什么辣不辣。

不辣。

陈敏回去煮面。煮的时候回头看了一眼。两个人在讨论什么。夹克男说了一句话，衬衫男还是没回应。只是看着对岸。

面端上去了。两个人在石墩子上吃。吃得很快。不像在品尝食物，像在完成一个步骤。吃完了，夹克男放下碗。碗里还有半碗汤。他说，走吧。

他们走了。

陈敏收了碗。看着他们的背影消失在土路的拐角。一路留下了三个鞋印。皮鞋的鞋印很深。土路被上次的雨泡过，干了以后表面是软的。踩上去是一个清晰的印子。


\scene

傍晚，太阳快要完全落下去的时候，方远来了。

他已经来了很多次。每次都是傍晚。他好像知道陈敏不忙了。每次来的时候店里都没有别人。他坐下来，点一碗面。吃得很安静。

今天也差不多。但今天他坐下的时候没有点面。先点了一支烟。陈敏没看他。他在面馆里只点素面，加辣。陈敏记住了。开始煮面。

他不会做饭。陈敏说。

方远把手里的烟弹了弹灰。烟在风里歪了一个方向。

你怎么知道。

女人都知道。方远朝厨房那边看了一眼。厨房其实不是厨房。就是锅和灶。碗摞在筐子里。

她会回来吗。

不知道。方远说。

陈敏没有追问。她只是把面端上去。方远接过去的时候，手指碰到了她的手指。是凉的。气温不低，但他的手指凉得像是从冷水里捞出来的。

面吃了一半。方远看着江面，问陈敏。今天有人来吗。

有个男的。穿白衬衫，皮鞋。在江边站了很久。还有一个跟班的人。看开发区，看地图。

方远把筷子放在碗上。筷子的一端搭在碗沿。

陈敏说，像投资人。看了一会儿就走了。

方远没有追问。只是把筷子拿起来，继续吃面。

天快黑透了。河面上只剩下最后一道光，像是有人用刀在水面上划了一道口子。陈敏去关卷帘门的时候，看见一个人从江边走过来。脚步声很重。不是方远的。方远的脚步很轻。

她回头看。

是一个中年人。穿着深色夹克，手里提着一个公文包。公文包看上去不新，角都磨白了。走过来的时候脸上有一层汗。天不热。大概走了一段。

他站在面馆门口。看了看关到一半的卷帘门。

还有面吗。他问。

陈敏看了他一眼。没什么特别的。跟其他来吃面的人差不多。年纪比方远大，比老赵可能还大一点。也是黑黑瘦瘦。头发短。眼睛很亮，不是那种脑子转得快的亮。是那种看东西看清楚了的亮。

什么面。

素面有吗。

有。要不要辣。

不要。

陈敏把卷帘门推上去一半。进去。灶上还有点热水。她把面下了。动作比下午慢了。一天下来，手有点酸。

面端上来。那人坐在方远的斜对面。方远还在吃，没抬头。那人把公文包放在脚边，拿筷子。

陈敏注意到他吃东西的方式。很仔细。每一口都嚼够了再吞。不像肚子饿，像是在品味某种他很少有工夫品味的东西。

吃了几口，他抬起头。

面不错。

陈敏说，水好。

他点了点头。没再说话。继续吃。

这时候方远已经吃完了。他站起来，走到陈敏的灶台前面，把钱放下。钱叠得很整齐。不是方远叠的，是银行出钱时的状态。他没有零钱。

他往他的面包车走。车停在江边。开锁，上车。发动机发动了。车灯亮起来。车灯照在江面上，照出一小片水面。那片水面上飘着一些浮萍。

车灯又灭了。发动机也熄了。

方远走回来。站在面馆外面。

陈敏抬头看他。

方远没说话。他只是看着那个正在吃面的人。看了一会儿。

方远说，辣不辣。

不辣。那人回答。

方远对陈敏说，给他加点辣。算我的。

陈敏愣了一下。拿了一罐辣椒酱，放在桌上。那人看了方远一眼，也看了一下辣椒酱。拿了一勺，浇在面上。搅了两下，吃了一口。

辣。

辣的比不辣的好吃。方远说。

那人点了点头。

他继续吃。吃得流汗。用手背擦了擦额头。汗擦到袖子上了。

一碗面吃完。他站起来。问多少钱。

方远说，陈姐的面，不贵。收他十块。

那人从口袋里掏钱。摸了一圈，没有零钱。拿出一张一百的。陈敏找给他九十。

方远看着。等陈敏找完钱，他才开口。

你做什么的。

那人说，我叫周立。我是林氏的人。

方远没有说话。陈敏洗了碗，碗放在灶台上靠锅的位置。江风吹过来。水面上已经什么都看不见了。她不认识什么林氏。但她知道这个名字。前两年江边挂牌子的那几个地皮上，旁边竖的广告牌写的就是这两个字。

方远说，你来这里做什么。

周立说，看地。

方远没有说话。他走回车边。这次把车发动了。

陈敏看着他开走。

车尾的红色灯渐渐变小，淡进夜色里。周立也走了。公文包还放在面馆门口。他又走回来，拿了包。说谢谢。说面很好吃。说要是有的话下次还来吃。

陈敏没应。只是看着他走进黑暗中。她把卷帘门拉到底，锁了。灶台擦干净。碗摞好。关灯。

她站在面馆门口。白天看到的那段江面上什么都没有了。渔船都收了。灯光都熄了。只有江还是那条江。黑沉沉的。看不出来是静止的还是在流。

但江一定在流。她知道。

她锁了门。沿着土路往街上走。街上有路灯。橙色的。她的影子在路灯下从一个变成两个，又从一个变成无数个。直到进了巷子，影子彻底消失。

\clearpage

% ==== Chapter 7 ====
\chapter{赌注}
\indent
转运中心的位置是方远自己选的。

在下游六十公里。一个叫梅镇的地方。说是镇，比村子大一点。过去有煤矿，煤挖完了，矿工走了，镇上的人走了一半。沿江那一段有三座空仓库。红砖盖的，房顶是石棉瓦，窗户的玻璃碎了大半。

方远站在最大的一座仓库前面。门是铁的，锈了一半。他用脚踢了一下，门开了。里面有一阵发霉的味道。墙壁上长着青苔。地上的水泥碎成拼图，踩上去不平。

鲁建国站在他旁边。看着这个破仓库。他什么都没说。

方远走进去。抬头。房顶有一个洞。洞把外面天空的一块漏进来。那一块很亮。周围黑。

多少钱。方远问。

一个当地的中间人站在门口。说，三座仓库，一年的租金十万。

方远说，买呢。

中间人愣了一下。买。

方远说，三座仓库，加前面那片空地。总共多少钱。

中间人走进去打了一个电话。打了好久。出来的时候脸上有点汗。

他说，空地不单独卖。如果要买，三座仓库加空地，八十五万。

方远说，七十五。

中间人又打电话。这次短了一些。出来说，八十万。最低了。方远没有说话。他看着那一片空地。空地的尽头是江。岸边有水泥的码头。是旧的，裂缝不比渡口那个少。但位置好。水够深。能停大船。

签合同。方远说。

回去的面包车上，鲁建国一路没说话。车开了二十公里，他才开口。八十万。方远，我们账上还剩多少。

不够。

那怎么办。

方远说，借。

鲁建国没有问找谁借。他知道方远一定有办法。方远没有说。只是开车。


\scene

一个月后，钱凑齐了。借了三个人的钱。一个是他父亲在山东的一个老战友，打了二十万。一个是老赵，打了十五万。一个是周学文，没说借什么理由，只说江湖救急，打了十万。方远自己的积蓄，加上鲁建国投进去的，正好八十万。

方远站在那座仓库前面。合同签了。钥匙在他手里。钥匙是铁的，很旧，上面有铜锈。他把钥匙揣进兜里。

第一天，他和鲁建国在里面待到了天黑。修好了门，换了锁，把碎玻璃扫掉。鲁建国去买了新石棉瓦，补完了那个屋顶洞。站在地面往上看，补丁是新的，颜色比周围淡，像一个刚结的疤。

转运中心开业的那天没有仪式。没有花篮，没有红布。只是开门。

来的第一批货是老赵的。六吨化工原料。从渡口出发，方远开着船，沿江往下走了六十公里。在梅镇的码头停靠。他和鲁建国两个人卸了货。六吨。三百袋。每袋四十斤。搬了三个小时。搬完的时候手抖得拿不稳水杯。两个人坐在地上喝水。鲁建国的嘴唇上全是灰。

他吐了一口唾沫，唾沫也是白的。

方远没有说话。他看着窗外。那个新补的石棉瓦在夕阳里反着光。


\scene

林氏的消息是周立传过来的。不是什么正式的通知。就是周立又来面馆吃了一次面。陈敏端面上来的时候，周立说了一句话。他说，我们公司要在下游建转运中心。

陈敏问，哪个下游。梅镇再往下十公里。

陈敏没有接话。只是在方远来吃面的时候，把这句话告诉了他。

方远听完了。吃了一筷子面。嚼完。

他说，十公里远。

陈敏说，近。

不远也不近。方远说，比我想的远了一点。

他继续吃面。吃完付了钱。站起来的时候陈敏看见他的脸。没有表情。


\scene

林氏的转运中心开业在一个月后。两条船。一条千吨级的，一条五百吨的。停在水面上，像两座小山。方远那天坐在自己的码头上，看着林氏的船从上游开下来。吃水很深，船头推开的水浪涌过来，摇动了他停在那里的旧船。

鲁建国站在旁边。也看了。他的表情没有方远那么沉。他的嘴唇动了动，没有声音。

第二天，老赵打来电话。说林氏的人找过他。价格低方远两成。

方远说，你怎么说。

老赵说，我说我的货只走你的船。

方远说，为什么。

老赵说，你那个疤能修多久。

方远说，没钱赚的事了。老赵说，那时候你不修房子上的疤。现在我不换。

方远挂了电话。鲁建国在旁边听着。听到了最后一句话。他没有问什么意思。但眼睛里有点湿。


\scene

林氏的动作很快。第一个月，方远丢了两个客户。水泥厂和纸箱厂都转了。电话通知的。电话很短。一个说，下次再合作。一个说，对不起。

方远没说什么。放下电话。在桌上的账本上画了两条横线。把这两个客户的名字划掉了。账本上现在只剩老赵的名字和一个新的名字：邻市港运。

邻市港运是周学文介绍的业务。每个月五趟。运量不大。但固定。不会跑。

鲁建国说，我们不能降价吗。

方远说，降不起。他们有两条船。

鲁建国说，那怎么办。

方远没有说话。他在桌子上铺开一张纸。不是地图，是一张时间表。列出了每一条航线的出发时间、到达时间、装卸时间、返回时间。他用笔在几个时间上圈了一下。

凌晨两点到四点。他说。

什么。

这个时间段，没人跑。

鲁建国看着时间表。看了一会儿。明白了。


\scene

方远决定不降价格。不改航线。只改时间。

他从第二天凌晨开始接货。凌晨三点。方远一个人开着船，从梅镇的码头出发，往上游走。夜间江上没有灯。只有他船头的那一盏。光照的距离有限。照不到的地方全是黑的。江水的声音比白天大。不是水变急了，是没有别的声音压住它。

他开得很慢。方向盘握得很紧。眼睛盯着前方。盯久了，光里面的水就会变得模糊。然后他把眼睛移开，看一会儿黑。再看回来。水又清楚了。

凌晨四点，船到江城。货主的人已经在码头等着。几个人。靠着车，缩着脖子。凌晨的风从江面上吹过来，比市区冷三四度。

方远靠岸。系好缆绳。跳上码头。工人开始卸货。他在站旁边看着。身上只有一件薄外套。凌晨的风吹得他的头发全乱了。他不理。

货卸完了。工头问他，明天还是三点。

方远说，三点。

工头说，你每天都这样。

方远说，嗯。

工头看着他。想说什么。没说。只是点了点头。上车。车开走了。尾灯像两粒红色的豆子。

方远开着船回梅镇。回程是一个小时。到了天已经亮了。江面上从黑变灰，从灰变白。那种白是一层一层的。先是最远的地方白一点，然后慢慢地，白往近处走。等到他靠岸的时候，整个江都是白的。

鲁建国在码头上等他。给他递了一碗面。

哪里来的面。

自己煮的。鲁建国说。

方远接过来。吃。面很软。煮过了。但他什么都没说。把面都吃完了。汤也喝了。碗底有一点沉淀，是碎的葱花和碎的面渣。

他说，还行。

鲁建国说，我从陈敏的面馆学的。

方远看了他一眼。看了一眼，没说话。


\scene

季节已经走到夏天了。五月。天亮得早。方远十二点就醒。他配货排线清账都在下午做。晚上继续开船。

一个月，他瘦了十斤。鲁建国更瘦。鲁建国瘦得脸颊凹进去。以前有一个下巴。现在下巴往上翘。像被两只手托着。

有一天他穿着雨衣回梅镇的仓库。进了门，倒在地上，说不在乎吃饭。鲁建国把他拖到板床上。过来摸他的额头。不烫。鲁建国拿筷子把他的嘴撬开。用茶缸把温水倒下来。喝完，慢慢睁开眼。说没事。两点还要接货。

你不去。鲁建国说。

方远说，你去。

鲁建国不会开船。

方远说，你学。

鲁建国看着他。看了很久。他不是第一次看到方远眼睛里那种不是恳求而是信任的表情。把茶缸放下了。问他还能走吗。方远站起来。走到船边。指着方向盘和仪表。指了半个小时。

鲁建国听着。记住了每一个东西的位置。油门在哪。转向在哪。灯怎么开。万一熄火怎么办。他听完了。说，我来。

那天晚上两趟都是鲁建国跑的。方远躺在板床上。听着柴油机的声音从码头上传来，然后远了。远了以后就只剩下江水的声音。他闭着眼睛。没有睡着。他知道鲁建国会开回来。一定。

天还没亮的时候发动机的声音回来了。由远而近。节奏很稳。方远从床上坐起来。看着鲁建国把船靠得漂漂亮亮。他们只是点了点头。


\scene

六月，第三个客户回来了。

是纸箱厂的。电话打过来，是一个女的。说话很犹豫。说想回来，问能不能把上次把给你们的部分让回来。方远说原来的价格不变。女的愣了一下。说林氏给了更低的价格。方远说我知道。她问那你为什么不做。方远说，因为你的货不止运一次。你要保证每次都能运到。你在乎价格，也在乎船会不会迟到。货站在雨里等两个小时，损失的比价格多得多。那边沉默了很久。说，好。

方远在自己的本子上写下了纸箱厂的新批次。

七月，第四个客户是一个新的名字。粮油中转站。他们是自己找上门的。来的是一个胖的男人。大肚子，穿花短袖。握手的时候手上全是汗。他说，我们是听一个工头说的。说你们夜里也跑，下雨天也跑，从没迟到。方远说那行。那你来什么时间方便。他说早上十二点。方远说十二点没有船。只有两个时间。胖男人说三点。方远说三点。定了。

八月，第五个客户也签了。

鲁建国数了一下本子上的名字。五个客户。超过了雨前。

他说，林氏没有再降价。方远说，他们还在观望。鲁建国说观望什么。方远说，观望我们还能撑多久。

鲁建国说，我们能撑多久。

方远说，说了不算。做下去。才能知道。

那天傍晚方远又在码头上洗船。船身上有淤泥，是从江底刮来的。他用刷子刷了半个小时。刷到天快黑了才停。鲁建国从仓库出来。把他拉到面馆去。

方远坐在面馆的老位置。陈敏煮面。端上来是两碗大排面。

方远抬起头看她。

陈敏说，多的这碗我请鲁建国。

鲁建国说，我呢。

陈敏说，你吃素面。

鲁建国笑了。方远也笑了一下。很短，但那个笑是可以被看到的。

鲁建国吃着素面。说，陈姐你知道这几天晚上他几点睡吗。

陈敏问他几点。

不睡。他说。我接货看上他站在码头上看水文。我开船回来，他还是站在码头上，看构图，看水位。我问他你是疯了吗。他说他不看充足就找不到路线。

陈敏看了一眼方远。方远低着头吃面。下巴有点胡茬。前几天刮过了，这几天又长出来了。

她说，想家了。

方远说，没有。

陈敏看了一眼江面。那天晚上天气很好。云很多。月亮在云的缝隙里露出一点。很快就又被遮住了。

鲁建国在抽烟。他说，陈姐你觉得我们能赢吗。

陈敏洗干净手。把指甲里的灰冲掉了。说，我说不好。但你们做的事情，跟别人不一样。

鲁建国说，什么不一样。

陈敏想了想。菜市场的菜，谁家的新鲜，谁家的不新鲜，吃一口就知道。一条街都卖菜，当天的最好卖。你们就是当天的菜。

鲁建国笑了笑。方远吃下最后一口面。拿起纸巾。

他说，谢了陈姐。

陈敏说，回去吧。早点睡觉。方远走的时候她在他袋子里塞了一张纸巾。方远捏了一下，没有回头。


\scene

八月下旬，方远的本子上又多出一个名字：林氏。

是周立打来的电话。他说想谈合作。方远约在面馆见面。周立来了。穿着林氏的公司衫。前面印着林氏集团四个蓝色字。他说他们有意转让几条下游的航线。问方远有没有兴趣。

方远听完。拿起桌上的茶喝了一口。茶烫了。

他说，你们不是来抢市场的。你们是来了解我的。

周立没否认。他笑了笑。这个笑和第一次见面时不太一样。里面有一点和善。也有一点害怕。

方远说完。放下茶杯。站起来。他走到面馆门口。叫住陈敏。陈敏过来。

他说，这个周先生的账，我给。以后他来都一样。吃面不要他钱。只要他来，就煮。

周立坐在桌上。看着方远的背影走进黑暗中。鲁建国看了一眼陈敏。陈敏没说话。把灶台再擦一遍。虽然已经很干净了。

\clearpage

% ==== Chapter 8 ====
\chapter{退路}
\indent
收购的消息是九月初来的。

鲁建国在码头上接到一个电话。他接的时候手上有油，是刚才检查发动机沾的。他把手机夹在耳朵和肩膀之间。听了几句，油从手上滴到手机上。

他挂了电话。走到铁皮房子门口。

方远在里面算账。桌上的本子摊开着，笔夹在中间。

刚才谁打电话。方远没抬头。

林氏。鲁建国说。

方远的笔停了一下。继续写。

他们说想收购。

多少。

鲁建国说了一个数字。

方远的笔停了。不是那种铅笔写完停的顿挫。是彻底停了。他把笔放在本子上。笔滚了一下，掉到地上。他没有捡。

鲁建国说，比市场价高五成。

方远没有说话。他弯下腰把笔捡起来。笔尖摔歪了。他捏了一下，没正过来。把笔放进桌上的笔筒里。笔筒是一个截断的塑料瓶，里面有三支笔。两支是黑色，一支是蓝色。摔歪的是蓝色那支。

鲁建国看着他。他知道方远在算。不是算钱。是算别的东西。方远在想的时候不会说话，不会动。像一个关掉了开关的机器。

过了大概三分钟。

方远说，再等等。

鲁建国说，等什么。

方远说，等他们再来一次。

为什么。

同一个价格不会只报一次。方远说，第一次是探路。看你的反应。你马上答应，他们会觉得报低了。你不答应，他们会加一点再报。第三次才是真的。

鲁建国说，你怎么知道是第三次。

方远说，不知道。猜的。

鲁建国看了他一眼。方远继续算账。算的是十月份的运输量。笔换了一支黑色的。


\scene

林氏的第二次报价比第一次高了一成。

这次不是电话。是周立亲自来的。开着林氏的车。车是黑色的，洗得很干净。停在码头旁边的土路上，轮胎上糊了一层泥。

周立走进铁皮房子。坐在那把椅子上。椅子还是鲁建国那年买的，座垫已经塌下去了。周立坐上去的时候往下陷了一下。他调了一下身体。没抱怨。

他说，这是我们最新的报价。

方远拿过来看了一眼。然后把纸还给周立。

他说，这个价格很合理。

周立说，不只是合理。方远说，但我不接。

周立看着他。脸上的表情很复杂。像是预料到了，又像是不愿意被预料到。他说，能问为什么吗。方远说，你们的条件里有一条。要求我退出这个行业至少三年。我不能签。

三年后呢。周立说。

方远看着他。三年后，这个行业还是这个行业。江还是这条江。但谁来运，不好说。

周立站了起来。他说，我会把这句话带回去。

方远送他到门口。周立上车前回头看了一眼那个铁皮房子。看了一眼，没说什么。车发动了。轮胎在泥里打了一个滑，稳住，开走了。

鲁建国走过来。他说，你刚才那句话是什么意思。

方远说，哪句。

谁来运，不好说。

方远说，字面意思。

鲁建国看着他。他知道方远不会再多说了。自己一个人去洗船。刷子在船身上刮出很大的声音。


\scene

九月的第二个星期，方远去了林氏总部。

没有约。他是直接去的。林氏的总部在省城开发区。一栋七层的楼。楼的立面是玻璃和铝板。门口有保安。保安让他登记。他在来访事由那一栏写了找林总。

保安打了一个电话。挂了电话说，林总没有预约不见。

方远说，你跟他说我是方远。你告诉他收购的事先别动。今天晚上我来找他喝茶。

保安又打电话。说了几句话。挂了。说，林总今晚有空。

方远点了点头。走了。他在省城的一家旅馆里住下来。旅馆很小。房间在二楼。窗户对着后院。后院里晒着几条床单。风吹过来，床单鼓起来，像帆。

天黑的时候，他把鲁建国给他熨的那件白衬衫穿上。在洗手间里蘸了点水，把头发顺到后面。镜子里的自己，他看了一会儿。不是在看好看不好看。是在看像不像一个能跟林氏老板喝茶的人。

他对着镜子系扣子。最上面那颗扣子，系了两次。


\scene

林老板住在开发区最里面的一栋别墅里。别墅不高，只有两层。但院子很大。院子的地上铺的是青砖，砖缝里有苔。院子里有假山。假山上面有水流下来。不是很大，就是一条。很细。流到下面的池子里，池子里有几尾锦鲤。

佣人带他进去。换了鞋。客厅很大。顶上吊着一盏水晶灯。灯亮着。方远看了一眼，没多看。他跟着佣人穿过客厅，走到后面。

茶室。

茶室不大。四面墙是木头的。颜色很深。放着茶叶罐、茶壶、一小只盆景。盆景是松树。很老了，树干歪着，像一个被时间压弯的人。

林老板在茶台后面坐着。

六十多岁。头发白了一半。脸上的皮肤松了，但精神很紧。穿着家常的深色上衣。袖口没有扣。不是不讲究，是讲究到了一定程度以后无所谓了。

坐。林老板说。

方远坐下来。茶台很矮。他坐下去的时候膝盖碰到茶台的边。他把膝盖移了一下。

林老板开始泡茶。泡茶的手法很慢。水先倒进壶里，热了壶再倒掉。再放茶叶。茶叶舒展开，像在呼吸。

方远看着他泡。没有帮忙。也没有说话。

林老板把第一杯茶推到方远手边。他自己拿了第二杯。两个人同时喝。茶是红茶。不浓。有一种隐约的甜。

好茶。方远说。

林老板点了点头。

沉默。大概有半分钟。

然后林老板说，你这个人。我把钱堆到你面前。你不动。我叫人去谈，你拒绝。我叫人加价，你还是拒绝。你到底想要什么。

方远放下茶杯。茶杯和茶台接触的时候没有声音。他说，我不是不要。是现在还没到时候。

还没到什么时候。

方远说，你们林氏的主业是地产。水运你们不做，也不懂。但下游在开发，政府对水运有要求。你们必须有一个水运公司在这里，哪怕只是做个样子。你们现在来买，不是因为你们想做这个生意。是因为你们必须让政府看到你们在做这个生意。

林老板把茶杯放在桌上。不是放，是顿了一下。杯底碰到桌面，声音不大，但有一点回音。

方远继续说，如果我把我的公司卖给你们。两年以后你们把地拿完了，水运这个板块你们没人管。许可证吊销，船卖掉，人走光。到那个时候，我拿到了一大笔钱。但我什么都没了。

那你今天来，是想告诉我你不卖。林老板说。

方远说，不是。我想告诉你，我不是来卖公司的。我是来卖位置的。

林老板看着他。看了很久。外面的水声还在响。很轻。轻到你不注意就听不见。

什么位置。

你们的水运负责人。方远说。

林老板的手指在茶杯边上停住了。停留了很长时间。长到方远数到了第四秒。

他拿起茶壶给自己倒了第二杯茶。也顺手给方远倒了第二杯。他没有说好还是不好。他说，你继续说。

方远说，林氏不是只有一个下游的开发项目。省城外面还有两个。别的地方还会有。每一个都需要水运。你现在缺的不是一个码头。是一个懂水运的人。

林老板喝了一口茶。把杯子放下来。说了一句。

你知道长江上有多少航运公司吗。方远说，大大小小上千家。

那你为什么行。

方远沉默了好一会儿。说，因为我不是只把水运当生意。对我来说，江是唯一的路。我没有别的地，没有别的资产。我只有江。你们的竞争对手也都有江。但你们的竞争对手跟江的关系，是租的。我的关系，是长在江上的。

林老板不再问了。他又泡第三泡茶。第三泡更淡。颜色从琥珀变成了金黄。

他们喝了第三泡茶。这时候已经十一点了。窗外的水声完全停了。假山的水应该是定时关的。

林老板说，那你说的位置。要我给什么。

一个平台。方远说。

什么样的平台。

林氏旗下的水运公司。独立核算。有人权财权。我负责。我不要你的收购款。我要经营权。方远把话说得很慢。每个字都单独吐出来。

林老板给他倒了茶。第四泡茶几乎已经没有颜色了。但他还是倒了一杯。

他说，你知道上一个跟我提这种条件的人现在在哪里吗。

不知道。

在拘留所。林老板笑了。方远也跟着笑了一下。他低下头，拿过茶杯，喝了一口淡得几乎感觉不到味道的水。


\scene

方远从林老板家出来的时候已经快十二点了。

他没有叫车。走路。省城的晚上和江城完全是两个样子。江城晚上除了江声什么都不剩。省城晚上还有霓虹灯、夜宵摊、酒吧门口排队的人。但他都看不见。他的脑子里只有那间茶室的声音。木头墙壁的颜色，茶香的味道，林老板的笑。

走到旅馆的时候，他在楼下的台阶上坐了一会儿。手机震了一下。鲁建国。

怎么样。

喝茶。

然后呢。

方远说，不知道。

鲁建国沉默了一下。你这个人，从来都说不知道。上次你说去省里批文，也说不知道。上次你跟林氏竞争也说不知道。上次你说纸箱厂会回来也说不知道。你每次说不知道，最后都是知道。

方远没有说话。

鲁建国说，算了，早点睡。电话挂了。

方远坐在台阶上。抬起了头。九月的天，云薄到透明。上面全是星星。那颗最亮的，像一粒钉子钉在天上。


\scene

第二天上午，他从省城回到渡口的时候，九月最大的那一批货还有两个小时到。他站在码头上。铁皮房子后面烟囱没有了，上一场台风把烟囱刮飞了。台风一过，秋天就来了。

空气里没有雾。江面清朗。对面码头的蓝色吊臂在动。

国庆节过后，林氏的第三个报价来了。不是电话。也不是周立。是一封邮件。发到方远的手机上。标题空着。正文很短：速联系。

方远读完了。把手机放下来。鲁建国在旁边抽烟。他问说，第三个报价了。你接不接。

不接。

那你怎么说的。

方远说，我上次就跟他们老板说清楚了。

他说了什么。

方远拿了茶缸去接水。烧开了。热气从茶缸里喷出来。

说不接就是不接。

鲁建国不问了。他把烟掐了。烟掐在码头的水泥地上，留下了一个黑印子。


\scene

当天晚上，方远给林老板发了一条信息。很短。十个字。

下次来江城。我请您看江。

发出去以后，他把手机放到枕头底下，闭眼睛。江水的声音从江面上传过来。用旧了的耳朵已经不在乎那条线的距离了。

他很快就睡着了。

梦中他站在一块陌生的码头上。水涨了三米。但他站的这个码头，一点都没被淹。

\clearpage

% ==== Chapter 9 ====
\chapter{转让}
\indent
十月。江上的雾比上个月浓了。

方远站在码头上。他身后是梅镇的三座仓库。最大的那座已经修好了。换过房顶，补了墙，装了灯。灯是白炽灯，照出来发黄。但比黑好。仓库的门口挂了一块木板。木板上用黑漆写了三个字：方记码头。

字是鲁建国写的。写得不怎么样。码字歪了，头字小了。但方远没有让他重写。

转让的决定是他一个人做的。


\scene

十月十五号。方远把鲁建国叫到仓库里。

仓库里有两把椅子。一把是他坐的，椅腿上缠着铁丝。一把是鲁建国的，座垫塌了。两个人面对面坐着。中间是一个茶缸。茶缸上印着红双喜。已经掉了一半漆。

方远说，我要把梅镇的业务转让给林氏。

鲁建国没有说话。他看着方远。看了大概有十秒。然后把头低下去。看着自己的手。手上的裂口好了很多。不像刚来的时候那样全是口子。但指节变粗了。是搬货搬的。骨节那里鼓出来，像长了一个包。

他说，你再说一遍。

方远说，梅镇三条航线。老赵的化工，纸箱厂，粮油中转。全部转给林氏。

鲁建国说，这是我们最好的业务。

嗯。

你要转给他们。

嗯。

鲁建国站起来。他走到仓库门口。站了一会儿。走了回来。又站了一会儿。又走回来。他蹲下去。蹲在椅子旁边。两只手抱着膝盖。指节上的骨包顶着裤管。

他问为什么。

方远说，因为这些业务赚的是辛苦钱。每趟几百块。一个月跑下来，扣完油钱和人工，剩不下什么。

鲁建国说，但我们做得好好的。客户回来了。业务在涨。

方远说，做得好，是因为没有人跟我们抢。现在林氏来了。他们有钱，有船，有人。我们做得再好，利润也会被压下去。最后变成谁价格低谁活。我不玩这个游戏。

鲁建国看着他。

那我们做什么。

方远说，做码头。

码头不就是在做吗。

方远摇头。他说，我们现在做的不是码头。是运输。运输是跑腿的。码头是坐着的。跑腿的永远比坐着的累。坐着的赚得比跑腿的多。

鲁建国想说点什么。嘴张开又合上了。他脑子里有很多话，但每一句拿出来都不对。他跟了方远快一年了。从那个雨夜开车来到渡口算起。这一年里他搬货、开船、修码头、接客户。手上起了茧，腿上落了伤。他以为他们在建一个运输公司。结果方远说，不做了。

他看着方远的脸。方远的脸很平静。像江面。不管水底下多急，面上是平的。鲁建国在这张脸上看不到犹豫。也看不到兴奋。只有一种他已经很熟悉的东西。确定。

鲁建国问，那我们做什么。

方远说，做中转。

方远没有说是，也没有说不是。他站起来，走到仓库门口。外面的雾散了一些。江面露出来。水很平。像一块灰色的布。

他说，你信我吗。

鲁建国说，信。

那就行。


\scene

方远给周立打了电话。说，可以谈了。

周立来的那天是十月二十号。他开着一辆白色商务车。车上还有一个穿西服的人。是林氏的法务。姓吴。戴金丝眼镜。手里提着一个黑色公文包。包很鼓，里面应该是合同。

他们在梅镇的仓库里谈。方远泡了茶。茶是老赵送的。不值钱，但能喝。

方远说，我转让三条航线的经营权。包括老赵的化工、纸箱厂、粮油中转。所有客户关系移交给林氏。

周立说，价格呢。

方远说了一个数字。比林氏的第三次报价低了三成。

吴律师抬起头。看了周立一眼。周立没有看他。周立看着方远。脸上是一种想不通的表情。他在想这个人为什么降价。

方远说，但有一个条件。

什么条件。

方远说，转让合同里加一条。梅镇码头的泊位使用权归我保留。林氏的船可以使用泊位。但泊位的经营权不转让。另外，我保留在码头建设仓储和中转设施的权力。

吴律师的笔在纸上停了一下。他看了一眼周立。

周立问，你留下泊位做什么。

方远说，做中转。

中转不是跟运输一样吗。

不一样。方远说，运输是谁跑得快谁赚钱。中转是谁的位置好谁赚钱。你们有船。但我有位置。没有码头，你们的船到了也卸不了货。

周立沉默了。他把茶杯端起来。没有喝。放在嘴边。放下。

他说，我需要请示。

方远点头。

吴律师走了。周立也走了。仓库里只剩方远一个人。他坐在椅子上。喝完了壶里最后一杯茶。茶凉了。有一点苦。


\scene

三天以后。周立打来电话。说林总同意了。但泊位条款需要改一下措辞。

方远说，措辞可以改。意思不能变。

周立说，你的意思是泊位使用权永久归你。

方远说，不是永久。是二十年。

二十年跟永久有什么区别。

方远说，有区别。二十年是一条线。线以内是我的。线以外的事，到时候再说。永久是不讲理。我不讲理。

周立笑了。他说，你这个人很怪。

方远说，我不怪。沿江六十公里，水深够位置好的泊位只有三个。你们有一个。港务局有一个。第三个在我这里。航线谁都可以跑。泊位跑不了。

周立没有再说话。他在电话那头沉默了。方远听到他翻纸的声音。可能是合同。可能是别的什么。

挂了电话，方远把鲁建国叫过来。

签了。他说。

鲁建国说，你真的要签。

方远说，合同里有一条。你注意看。是关于泊位使用权的。那一条比所有的航线都值钱。

鲁建国说，我不懂。

方远说，你不需要懂。你只需要知道一件事。以后我们的船不跑货了。我们做中转。别人把货运到我们的码头，我们卸货、分拣、仓储、再配送。赚的是位置的钱。不是力气的钱。

鲁建国想了想。

他说，你从什么时候开始想这个的。

方远说，从辞职那天。

鲁建国看着他。方远的脸很平静。像江面。不管水底下多急，面上是平的。


\scene

签字的那天是在陈敏的面馆。

方远选了这个地方。不是因为他喜欢面馆。是因为他觉得一个大的决定应该在一个跟它不匹配的地方做。太正式的地方会让人紧张。紧张的时候容易看错东西。

周立和吴律师来了。合同打印了两份。每一份有十二页。吴律师逐条念。念到第八条的时候，方远让他停下来。

第八条是关于泊位使用权的。吴律师念完。方远看了一遍。逐字逐句。

他拿出一支笔。在第八条下面加了一句话。写得很慢。字迹不好看，但很清楚。

他写的是：转让方保留在泊位周边五百米范围内建设仓储和中转设施的权利，受让方不得以任何理由阻拦。

吴律师看了。周立也看了。

周立说，这一条。

方远说，这一条是整份合同里最重要的一条。

周立想了一下。说，好。

两个人签了字。方远的签名很潦草。周立看了一眼，没认出来是什么字。但他没有问。

签完字。陈敏端了六碗面上来。素面。不加辣。周立吃了一口。说，好吃。

陈敏说，水好。


\scene

陈敏在擦桌子。桌子上有面汤的痕迹。她擦得很仔细。擦完一张换一张。抹布在桌面上画圆。画得很慢。她也在听。从头到尾。她什么都听见了。合同、泊位、转让。这些词她不太懂。但她听懂了一件事。方远把做了一年的东西交出去了。换回来的只是一张纸。

鲁建国吃了面。但没有说话。

签完以后大家走了。只剩鲁建国和方远在面馆里坐着。陈敏还在擦桌子。抹布画着圆。画得很慢。

鲁建国说，你把最好的业务给了别人。你留了一个泊位。我不明白。

方远说，你不需要明白。你只需要看。

看什么。

看明年春天。

鲁建国想追问。但方远已经站起来了。他把面碗端到陈敏的水池边。他自己洗了碗。

陈敏看了他一眼。她说，你什么时候学会洗碗的。

方远说，现在。

陈敏没有笑。她只是把抹布递给他。方远把碗擦干，放在碗架上。碗架是铁丝编的。碗放上去的时候铁丝弹了一下。

他说，陈姐，明年春天你再来看这个码头。会不一样的。

陈敏说，怎么个不一样法。

方远说，水还是那个水。但码头上的东西不一样了。

他没有解释更多。他走到门口。回头看了一眼面馆。四张桌子。折叠的。塑料贴皮翘了边。灶台上还有热气。陈敏站在热气后面。热气模糊了她的脸。但方远看得很清楚。

他转身走了。脚步声在土路上渐渐远了。

陈敏把卷帘门拉了一半。站在门口看了一会儿。看到方远的背影拐过了江边的弯道。看不到人了。她才把门拉到底。


\scene

那天晚上方远回到渡口。渡口的铁皮房子还在。码头的灯还亮。江水在灯下是暗黄色的。他把转让合同的副本从包里拿出来。翻到第八条。他加的那句话在页面的最下面。字迹很清楚。

他看着那句话。看了很久。

然后他把合同放进铁皮柜子里。锁上。钥匙放在衬衫口袋里。

他走出铁皮房子。站在码头上。风比白天冷。是深秋的风。从上游吹来。带着一种植物腐烂的味道。不是难闻。就是秋天江水的味道。

他蹲下去。手摸了一下江水。水很凉。不是刺骨的凉。是那种你摸了一下就会缩回来的凉。

他站起来。看了一眼梅镇的方向。什么都看不见。太远了。但他知道那个码头还在。三座仓库还在。泊位还在。

他把手机掏出来。给鲁建国发了一条信息。

明天开始，找施工队。

鲁建国很快回了一条。修什么。

仓库扩建。方远回。码头的东边还有两百米空地。全部盖仓库。盖两层。

鲁建国回了一个字。好。

方远把手机收起来。风更大了。他拉了一下衣领。往铁皮房子走。走了几步，停了一下。

回头看了一眼江面。

雾又升起来了。从水面上冒出来的。像一口气。先是薄薄的一层，然后变厚。把灯光挡在后面。灯光在雾里变成一团圆圆的晕。像月亮掉进了水里。

他进了门。关了灯。

黑暗里只有江水的声音。那声音一直在。从他第一天来这个渡口的时候就在。不管白天黑夜，不管涨水退水。江水发出的声音永远差不多。不紧不慢。像一个一直在说话但没人注意听的人。

方远躺在行军床上。眼睛睁着。黑暗里什么也看不见。但他知道窗外是江。江上有雾。雾后面是对岸。对岸的灯光比一年前多了。开发区在建。塔吊的灯是红色的，在夜空里一明一灭。

他闭上了眼睛。

\clearpage

% ==== Chapter 10 ====
\chapter{水声}
\indent
又一年。十月。

方远站在梅镇的码头上。码头已经不是一年前的样子了。

东边的两百米空地被填平。铺了水泥。盖了三排仓库。每排六间。一共十八间。灰色的铁皮墙，蓝色的屋顶。仓库之间隔着很宽的路。路够两辆货车并排走。路面上画了白线。白线很新。

码头上的吊车换了。以前那台是从老赵的厂里买的。旧的。吊臂转起来吱吱响。现在是两台。新的。电动的。没有声音。吊臂转起来像一只安静的手臂。

但他最喜欢的东西还是那些旧的。水泥码头上的裂缝填了，但缝的纹路还在。缆桩没换。锈还在。只是刷了一层防锈漆。红色的。在灰色的码头上很扎眼。

他每天早上还是五点起。比以前多了一项。到办公室先看报表。他的办公室在第三号仓库的二楼。房间不大。一张桌子，一把椅子，一台电脑，一个茶缸。墙上挂着一张图。是沿江六十公里的地图。图上有三个红色的点。梅镇码头、渡口、下游转运中心。三个点连起来像一个钝角三角形。他把这个三角叫做他的江。

窗外是江面。窗户很大。是他专门让人改的。以前的窗户只有一半大。他让工人把墙打了，换成整块的落地玻璃。阳光好的时候，整个办公室都是亮的。阴天的时候，江水的颜色映进来，像一面灰色的墙。

他坐在窗前。看着江。


\scene

一年前他把三条航线转给林氏。很多人看不懂。鲁建国看不懂。陈敏看不懂。老赵倒是没说什么。他只是问了一句，你还做不做化工。方远说不做了。老赵说好，那以后我的货找谁。方远说找林氏。老赵说行。

后来老赵跟他说了一件事。他说他有一次经过梅镇码头。看见码头上全是车。十几辆。排队等着装货。他那天要去别的地方，就没停下来。但他看了一眼。他说他看见了方远站在码头上。手里拿着一只对讲机。周围的工人跑来跑去。方远就站在那里，像一根桩子。

老赵说，你站着不动的时候，像一个指挥交通的警察。

方远笑了一下。说，我不指挥。我只是看。

现在的码头确实不需要指挥。人多了。鲁建国管调度。他手下有八个调度员。每个人负责一个区域。码头上的货分三种：散货、集装箱、冷链。散货是老本行。集装箱是后来加的。冷链是一笔意外的收获。对岸有一家水产公司。保鲜要求很高。全省只有三家码头能做冷链转运。方远是第三家。

他把码头做成了平台。不是他不做了。是他不亲自做了。他自己造船。造了一艘大船。三百吨。不用来运货。用来做浮动仓库。船停在泊位上。货卸到船上。分拣好了再装到小车。每天处理上百吨的货。油钱不便宜。但位置便宜。码头上每处理一吨货的净利润，是当初跑运输的五倍。

他把这个叫水的逻辑。鲁建国问他什么逻辑。他说水不走直线。水找阻力最小的路。

鲁建国说，那我就是你找到的阻力最小的路。

方远说，你不是。你是帮我修路的人。


\scene

十月底的一天。下午。阳光很好。江面上泛着白光。碎碎的。像一片一片的鱼鳞。

方远在办公室里看报表。电话响了。门卫。说门口有一个人要找方总。没有预约。穿着一件深色中式外套。头发花白。

方远说，让他进来。

他没有去码头接。他坐在办公室里。把茶缸里的茶倒了。换了一壶新茶。红茶。不太浓。有一种隐约的甜。

林老板进门的时候，方远站起来了。

林老板环顾了办公室一圈。没有说好。也没有说不好。他只是在窗户前面站了一会儿。看着外面的江。看了一分钟。也许更长。

他说，你把窗户开这么大，冬天不冷吗。

方远说，冷。但看得清楚。

林老板点了点头。坐下来。方远给他倒了茶。茶杯推过去。林老板拿起来。没有马上喝。他把杯子放在鼻子前面闻了一下。

红茶。

嗯。

上次在你这里喝的，也是红茶。

方远没有说话。林老板喝了一口。

他说，我今天不是来谈生意的。

嗯。

我就是想来看看。

方远说，那你的感觉呢。

林老板把杯子放下。看着窗外的江。

他说，一年前你来找我。说要把最好的业务转给我。那时候我以为你是撑不住了。想拿一笔钱走人。后来你加了一个条款。留了泊位。我后来让法务把条款查了一遍。法务说没问题。

嗯。

但我总觉得哪里不对。说不上来。今天看了你这个码头，我明白了。

方远说，明白什么。

你从来没想过要卖。你说转让，其实是把你不想做的业务甩给我。你自己留了最值钱的东西。码头。位置。基础设施。你在林氏之前就看到了。下游一开发，走水运的货会多。多出来十倍的货。但泊位不会多出来十个。泊位还是那三个。我有一个。港务局有一个。剩下的，在你手上。

方远没有说话。他把茶壶拿起来。给林老板的杯子里添了茶。林老板看着茶水从壶嘴里淌出来。水流很细。很稳。没有溅出来。

林老板说，你把转让的条款写得滴水不漏。泊位使用权归你，仓储权归你，周边五百米归你。我的人在码头上的船必须经过你同意才能停。你的水运公司已经不是你一年前跟我说的那个小码头了。

他停了一下。说，你现在是整个沿江六十三公里最好的中转港。

方远说，是靠你们的货。

林老板说，是大资本。他笑了。他的笑里有一点苦涩。是在笑自己。也是在笑方远。他说，我六十多岁的人了。被你四十岁的人算计了一遍。

方远说，我没有算计你。我算计的是水。

水？

嗯。方远说，水往低处流。谁的位置低，水就往谁那里流。我跟水学了两年。学到一件很简单的事。不要追水。等水流到你这里。

林老板沉默了很久。窗外的江水在响。很轻。林老板说，我没有见过你这样的人。

方远说，什么样的人。

林老板说，能等的人。


\scene

他们在办公室里坐了一个小时。喝了两壶茶。说了一些别的。林老板问了他的家人。方远说父亲很早就走了。跟母亲不联系。林老板没有追问。

林老板走的时候，方远送他到码头门口。车等在外面。黑色的。轮胎上沾了泥。林老板上车前回头看了一眼。不是看方远。是看码头上那台旧的吊车。红色的。缆桩上的防锈漆还没干。

他说，那台吊车是你最早买的那台。

方远说，是。

林老板说，为什么不换。

方远说，还能用。

林老板上了车。车窗摇下来。他说，下次来，请我吃碗面。这里有一家面馆，我以前吃过一次。

方远说，陈姐的面馆。

嗯。上次是你们签转让合同的时候。我记得她面里放的水好。

方远说，不是水好。是面好。

林老板笑了笑。车窗升上去。车开走了。轮胎在泥里打了一个滑。然后稳住了。尾灯消失在弯道后面。

林老板走了以后，方远一个人回到办公室。坐回窗边的椅子上。江水还在流。光变了。太阳往西斜了。江面从碎白变成了金黄。像镀了一层铜。

他给鲁建国发了一条信息。

林老板来过了。

鲁建国马上回了一条。他说了什么。

他说我算计了他。

你怎么说。

我说我算计的是水。

鲁建国回了一个竖大拇指的表情。方远没有回。


\scene

傍晚的时候，方远走去了陈敏的面馆。

面馆没有什么变化。还是那四张桌子。塑料贴皮翘边的地方多了一角。他用手指按了一下。贴皮又翘起来了。他没有再按。

陈敏在灶台后面。

一碗素面加辣。他说。

陈敏看了他一眼。说，你今天看起来很高兴。

方远说，是吗。

嗯。你笑的时候眉毛不动。但嘴角会往上走一点。现在你的嘴角在往上走。

方远摸了摸嘴角。没说话。

陈敏把面下锅。水开了。白色的蒸汽从锅口喷出来。她用筷子搅了一下。面条在水里散开，像一朵花。方远坐到了旁边的椅子上，看着窗外。窗外能看见一小截江。是弯道的尾巴。江水从那里流过，碰到岸边的石头，打出很小的水花。

面好了。端到他面前。碗还是那种蓝色的搪瓷碗。碗沿上有磕掉的瓷，露出下面的铁。方远端起来。吃了一大口。

他说，陈姐，你记不记得我第一次来是什么时候。

陈敏想了想。说，三年前。你辞职那天下午。点了素面加辣。

嗯。

你那时候什么都不说。我看你吃了一整碗面，一句话也没说。后来你走了，我在门口看你。你站在路灯下面。风吹过来，你没动。

嗯。

陈敏说，变了。你说话比以前多了一点。

方远说，面没变。还是那个味道。

陈敏说，面不会变。人变。

方远低着头吃完了面。把汤喝掉一半。抹了抹嘴。搁下筷子。他的手搁在桌上。手指很长。骨节很粗。指甲里有灰。是码头上的灰。洗了也还是有。

他说，陈姐。我想问你一件事。

什么。

你的面馆，以后有没有想过扩大。

陈敏擦了擦手。手上的抹布拧了一下。水从抹布里滴下来。

想过。但不知道怎么做。

方远说，码头的人越来越多了。每天吃饭的人排队排到门外。你现在一个人做不过来。再加一个人。前厅加几张桌子。灶台加一口锅。旁边那间铺面上一周刚退租。你可以盘下来。

陈敏看着他。她说，你帮我看房子啊。

方远说，不是帮你看。是帮你算。

陈敏转过头去。她的手在围裙上擦了又擦。围裙本来就不脏。她擦了这么久，是因为不知道把手放在哪里。

你不管我的面馆，我就不帮你算了。方远吃完了面。站起来了。从碗底下把筷子抽上来，摆好。二十块钱放在桌子角上。那个角上一直都有陈敏放的一小瓶醋。钱放在醋瓶子旁边。他没有把醋瓶碰倒。

他说，土豆牛肉盖饭。我出去的时候跟你说。你心里有数。

陈敏说，什么有数。

方远说，扩建的事。你心里有数。只是没有说出来。

他走了。门在身后吱吱响。

陈敏把碗收走了。碗里的汤还温的。


\scene

十一月初。天凉了。江上的雾又加了一层。

方远把码头上的事交给了鲁建国。自己去了一趟省城。不是去找林老板。是去省交通厅开一个会。关于内河水运标准化的研讨会。方远本来不是正式代表。但主办方打电话来，说想请方记码头的方总来分享一下经验。

方远没有准备PPT。他拿了一张纸。上面写了几个字。字不多。会场灯光暗，他眯着眼念：码头比船值钱。中转比运输值钱。泊位不要卖，留住。不要跟水较劲。等水流到你这里。

开会的时候，底下的人很安静。他说完以后没人鼓掌。他以为自己的发言太短了。太简单了。但散会的时候有好几个人拦住他。递名片。加微信。说想参观他的码头。

方远把名片装进口袋里。回到旅馆。他坐在床上。把名片一张一张拿出来看。有七张。三张是省内的。四张是外省的。外省来的。有一个是湖南的。有两个是江西的。还有一个是重庆的。

他把名片翻过来。反面是空白的。他拿起笔。在最上面那张名片上写了几个字：等水流到你这里。

他把名片压进了本子的夹层里。


\scene

十二月。鲁建国跟方远说，他要请两天假。方远说做什么。他说回去看看。他知道他说的是谁。他老婆。

去吧。方远说。开我的车去。路不好走。

两天以后他回来了。方远在办公室里等他。

他说，她问我你在做什么。

你怎么说。

我说你在做码头。她不信。

方远说，你把今年的流水单给她看了吗。

鲁建国说，没有。我只说了一句。方远答应我的，都做到了。

方远转过头来看着他。鲁建国站在门口。外套是新的。灰色的夹克。袖子有点长。他说是他老婆给他买的。她说赏的，买小了，给他是刚好。

方远说，她让你回来吗。

鲁建国摇了摇头。说，她在等我辞职。

方远没有说话。沉默了很久。然后说，不急。

鲁建国说，嗯，不急。


\scene

除夕那天，方远订了十五盒饺子。码头上有十五个没回去过年的工人。方远没回去。鲁建国也没有。两人在铁皮房子里吃饺子。饺子是韭菜鸡蛋馅的。鲁建国说是他老婆包的。走前放在冰箱里冷冻。让他除夕拿出来煮。

方远说，你老婆其实挺好的。

鲁建国说，嗯。只是不信我。

晚上十二点。远处放起了烟花。从江面上爆开来。红色绿色金色。在水里倒映成碎成一片的光。方远站起走到门口。仰起脸看了很久。烟灰从天上飘下来。他闻到了火药味。这个味道让他想起了小时候。他也说不清为什么。河水安静地托着倒影。

鲁建国也出来看着烟花。

他说，你今年过年一个人。

嗯。

你不想找个人吗。

方远看着烟花。他说，我找到了一条江。

鲁建国没有再接话。


\scene

三月初。江水解冻。水温回到七度以上。

方远的码头上多了一个招牌。方记码头有限公司。牌子上有工商注册号和经营范围。经营范围写得很长。水运、仓储、中转、装卸、报关。每一项后面都有一个对应的章。

有一天下午。陈敏走了过来。方远在码头上指挥装货。对讲机挂在脖子上。陈敏在他身后。看着他。看到他跟工人说的话不多。都是一两个字。卸、换、起、放。工人知道他的意思。只要他说一个字，它们就知道做什么。很精确。

方远转过身。看见陈敏。他说，吃饭吗。

陈敏说，我旁边的铺面盘下来了。

方远把对讲机从脖子上拿下来。

多大了。

十二张桌子。灶台三口。加了一个伙计。是以前对岸水泥厂食堂的。做盖浇饭很绝。

那你这就是开起来的第二家店了。

嗯。

方远说，那你的面馆有名字吗。

陈敏沉默了一下。说，有了。叫水还在流。

方远没有说话。他愣了一下。这个愣的时间很短。但足够一个人感觉到。陈敏感觉到了。但她没有问。

方远说，你什么时候想的这个名字。

陈敏说，你以前经常跟我说。你喝茶的时候说，吃面的时候说。你从来没有解释过这句话是什么意思。但每次你说的时候，你的眼睛都看着江。我就知道了。

方远没有说话。陈敏等他说话。等了大概一分钟。方远还是没有说。她把双手插进围裙兜里。转身走了。方远站在码头上。看着她的背影。围裙带子系得很紧。走到码头拐角的时候，她的身影被仓库挡掉了。

他低下头。对讲机响了。他按了一下。那边说2号泊位有三条船排队。问他放哪条先。

他说，先放最慢的那条。


\scene

一周后。方远又去了陈敏的新面馆。

新面馆的墙上挂着菜单的牌子。红色塑料板。白字。菜名写得很工整。饮料区下面还有一行小字。是陈敏自己写的：面好是因为水好。

方远看到这句话的时候笑了。他笑的时候眉毛还是不动。

站在他身后的一个工人端着一大碗面。他看到方远在自己对面坐下。问他，你以前是做什么的。很多人都问过他这个问题。他的答案永远一样。

跑船的。一年以后才做的码头。

工人哦了一声。以为他的老板不愿意多说。低下头去吃面了。

方远自己也点了一碗面。他要的牛肉面。陈敏端过来的时候说，旁边还有一个铺面也要转。她想盘下来。做成二十四小时的。因为码头上凌晨三点也有人要吃东西。夜班的装卸工。船工。货车司机。

方远说，能行。

陈敏说，你跟以前不一样了。以前你说能行的时候只有两个字。现在后面还加一句。

方远没有接这个话。他说，水还在流。

陈敏看了他一眼。她说，我已经卖了两千碗了。招牌菜是土豆牛肉盖饭和素面加辣。

方远说，素面加辣是你最早做的。

嗯。

陈敏说完这句话，转身进厨房去了。方远听见灶台的声音。火打着。油热了。菜下锅。滋的一声。像下雨。


\scene

十月底。一年整。

方远在办公室里算账。大船浮仓在码头上停满了一天。全天的中转量突破两百吨。破了记录。他把数字写在账本上的时候，笔没有停顿。写完以后把本子合上。放进抽屉里。抽屉里有一件旧东西。是那把磨掉了刻度的钢尺。他从规划局带出来的。尺面刮花了。上面的刻度一条都看不清了。但他没有扔。

他拿出钢尺。翻了一个面。背面有一行字，是写上去的，笔画不好看，但很清楚 ：水还在流。

下面落了日期。是他辞职那天下午四点半。

他把尺翻回来。尺面平滑。什么都没有。他拿起自己的手机。对着窗户拍了一张江水的照片。外面是灰白色的天。江水很平。没有风。

把照片发给了鲁建国。配了两个字：还在。

鲁建国很快回了一条。

我知道。


\scene

林老板是十二月份来的。第二次。

这次不是突然来的。他提前一天让人打了电话。方远知道他来。但没有安排什么特别的。码头正常运转。工人正常搬货。仓库的门开着。叉车在院子里跑来跑去。地面是湿的。刚冲过水。

林老板的车停在新修的停车场里。黑车。换了新的。比去年那辆大一圈。林老板下车的时候，方远在码头上的吊车旁边等他。

好热闹。林老板说。

每天这样。

林老板点了点头。他没有进办公室。方远带他在码头上走了一圈。从东面走到西面。看了三排仓库。看了浮仓。看了调度室。调度室是一间小房子。四面都是玻璃。鲁建国在里面。他看见方远带着一个人走过来。不认识。站了起来。有一点紧张。

方远说，这是林董。

鲁建国说，林董好。林老板点了点头。

方远然后带林老板走进了陈敏的面馆。十二张桌子。几乎坐满。方远找了一张靠窗的。两个人坐下来。陈敏走过来。烫的菜单都不用递。方远说两碗素面加辣。

林老板看了看墙上的字。

面好是因为水好。他说。

嗯。

是你的话。

陈敏写的。

林老板没有再问。他把筷子从筷筒里抽出来。用开水烫了一下。他的动作比方远多。先烫筷子，再烫碗，再烫小盘里的辣椒。每一道程序都做得很认真。方远看着。没有帮忙。

面来了。林老板吃了一口。他说，你这个面。

嗯。

跟你四年前请我那一次一样。没变。

嗯，面不会变。人变。

林老板放下筷子。他说，我跟公司的人说。水运这块不用再招人了。全交给你。你一年前怎么算的，我心里有数。别人说你抢了林氏的市场。我说他没抢。他建了一个新市场。然后把它让给了我。他在等更大的水流过来。

他用筷子挑了一撮面，慢慢吃完。你做的不是生意。你在改河道。这句话让方远沉默了几秒。

他说，我就是个跑船的。

林老板笑了笑。把碗里的面吃光了。汤也喝完了。他放下筷子。说，好。方远说，好。他们吃了第二碗面。这次两个人都没有说话。

外面。灶台下面炉子烧着。火很旺。


\scene

除夕。又一年。

陈敏的面馆开到很晚，深夜十二点才关门。码头上的人吃完年夜饭，陆陆续续走了。只剩方远和鲁建国。两个人在铁皮房子里烧了一壶茶。电炉上咕嘟咕嘟响。

很久没喝茶了。鲁建国说。

方远说，太忙了。

鲁建国说，去年忙也喝了。

方远说，去年喝的泡面。

鲁建国笑了一下。他拿起茶杯。茶很烫。所以他只是端着。没有喝。他说，你记不记得，那个下午我到渡口来找你。你一个人坐在码头上，坐在那半段水泥墩子上。我问你要怎么办。你说走走看看。我又问你看什么。你不说。

嗯。

后来有一个下午。在梅镇。你坐在码头上看江。我去找你。你才说了。你说在看水还在不在。

方远把茶拿起来。喝了一口。烫了一下舌头。他闭嘴。等了一会儿。说，水一直都在。

鲁建国说，嗯。水不会走。人会走。我把你也算是水了。

方远说，我不是水。

鲁建国拿着杯子碰了一下他的杯子。碰出来的声音很脆。

你是。

他放下杯子。看着窗外。窗外是江。江上有月光。月光碎在水面上。一明一暗。一明一暗。

方远没有说话。他把茶喝完了。站起来。走到外面。站在码头上。江水的声音从下面传上来。很轻。一直是这样。从他选择辞职的那个下午站到窗边看这条江开始，那声音从来就没变过。

他蹲下去。手伸进江水里。很凉。不是刺骨的凉。是那种感觉到凉就会把手缩回来的凉。他没有缩。他的手在水里多停了几秒。然后拿出来。把水滴甩干。水滴落到地上。只轻响了一下。

风停了。没有声音。

江还在流。

\clearpage

\end{document}
